\documentclass[11pt]{article}
\usepackage[margin=1in]{geometry}
\usepackage{amsmath,amssymb,amsthm,booktabs,tabularx,array,microtype,hyperref}
\hypersetup{colorlinks=true,linkcolor=blue,urlcolor=blue,pdftitle={Evidence, Influence, and Internal Capital Allocation},pdfauthor={Research draft prepared for Dr. Jan-Hendrik Weinert}}
\newtheorem{theorem}{Theorem}
\newtheorem{proposition}[theorem]{Proposition}
\newtheorem{corollary}[theorem]{Corollary}
\newcommand{\E}{\mathbb E}
\newcommand{\ind}{\mathbf 1}
\title{Evidence, Influence, and Internal Capital Allocation\\\large Rational Inference, Factual Verification, and Metric Governance}
\author{Dr. Jan-Hendrik Weinert}
\date{19 October 2026}

% Publication formatting edition, 19 October 2026. Scientific equations unchanged.
\usepackage{xurl,needspace,etoolbox}
\geometry{a4paper,margin=25.4mm}
\hypersetup{hidelinks,pdfauthor={Dr. Jan-Hendrik Weinert}}
\setlength{\parindent}{1.2em}
\setlength{\parskip}{3pt plus 1pt}
\setlength{\emergencystretch}{3em}
\hfuzz=0.2pt
\raggedbottom
\pagestyle{plain}
\makeatletter
\renewcommand{\@seccntformat}[1]{\csname the#1\endcsname\ifstrequal{#1}{section}{.}{}\quad}
\renewcommand{\maketitle}{%
  \begin{center}
  {\LARGE\bfseries\@title\par}
  \vspace{0.8em}
  {\normalsize\@author\par}
  \vspace{0.35em}
  {\small\@date\par}
  \end{center}
  \vspace{0.4em}}
\renewenvironment{abstract}{%
  \begin{center}\small\bfseries Abstract\end{center}
  \begin{list}{}{\setlength{\leftmargin}{0.5in}\setlength{\rightmargin}{0.5in}}
  \item[]\small\setlength{\parindent}{0pt}\ignorespaces
}{\end{list}\vspace{0.6em}}
\makeatother

\begin{document}
\maketitle
\begin{abstract}
Factual auditing can detect falsification while leaving accurate but selectively chosen metrics undetected. We solve Bayesian allocation with private quality and metric-selection costs, then optimize random auditing of positive evidence packets. The equilibrium distinguishes all-submit inflation, inflation/non-submission mixing and metric-only reporting. Submitted and passed masses differ, and failed checks can still lead to funding. We derive a finite global candidate set from regime boundaries and stationary equations. Convexity excludes an interior welfare maximum in the mixing regime, whereas an all-submit regime can have an optimal partial audit with residual inflation. The resulting information and resource frontier is compared with no evidence and locked metrics. Results are restricted to the stated evidence menu and uniform-cost technology; computations are synthetic.
\end{abstract}
\section{Introduction}\label{sec:intro}
Standardized investment evidence serves two different purposes. It can make a factual assertion reproducible, and it can determine which economic question is put to the capital allocator. Those purposes need not coincide. A number can accurately describe a selected customer group while being uninformative about a project's total value. The manager choosing the group may understand this perfectly, and the allocator may correctly anticipate the selection without being able to identify which individual packet contains it.

This article asks when verifying submitted facts improves a local internal investment decision. We separate falsifying a canonical measurement from selecting an accurate favorable measurement. The distinction matters because a factual audit detects the first action but, by construction, not the omitted measurement universe behind the second. The economic friction is incomplete observability of the evidence-generation recipe and of its private cost, not a fixed political preference in the allocator's objective.

The baseline comparison shows that full factual verification shifts strategic reporting to metric selection and raises real preparation costs without improving gross value. Every imperfect architecture in that binary, zero-shock menu is dominated by declining evidence. The full-range R2 solution permits valuable imperfect architectures, but establishes an important boundary: costless full factual audit cannot lower the value of an already valuable unaudited open architecture under uniform costs. That result remains intact. Positive audit fees do not rescue a disproved pure-substitution claim.

The source comparison now uses Wulf's June 2009 prepublication mechanism and the distinct fixed-false-claim coverage model of Ba\'nka and Chudziak (2026), with version and access limits stated in Section \ref{sec:literature}. The present main contribution optimizes a random factual-audit intensity within the same open-metric technology. We solve all sender and receiver regimes, including a previously omitted selection-versus-inflation cutoff, distinguish passed from submitted packets, and derive a finite global candidate characterization. The inflation/non-submission mixing frontier cannot have an interior maximum. An all-submit frontier can, and an exact-primitive witness gives a uniquely optimal partial intensity with positive inflation. A minimal deterrence probability is therefore not a sufficient governance recommendation. Neither is full factual checking merely because a restricted policy menu contained only zero and one.

This is a small one-candidate allocation game, not a firm-wide capital market. The new uniform-cost solution covers every finite outside-option shock width, including a separate zero-width derivation. The earlier distribution-class binary-audit and explicitly specified sponsor-signal results are retained within their stated scopes. Optimized open auditing is compared with the existing no-evidence and binary locked-metric policies. No transfers, penalties, arbitrary signal design or endogenous project quality are added.

The earlier manuscript's imposed score, deterministic sponsor premium, and ex ante diagnostic exercise are no longer the main model or supporting results. Their numerical examples, including the fixed-action triangular exercise, are not used as evidence for this revision. The present manuscript derives posterior beliefs, evidence actions and funding jointly after managers learn their information.

\section{Literature and the scope of the contribution}\label{sec:literature}
Wulf (2009) is the closest internal-capital-market predecessor. We now inspect the actual June 2009 author-paper text displayed in the public index, specifically Section 2, equations (6)--(10), and Appendix A (printed pages 37--40). Its cover explicitly identifies a prepublication version. This is substantive primary-paper reading through a third-party mirror, not authentication of that host's authorization, transcription, or identity with the publisher final. The July 2000 early working record, September 2005 SSRN working record, June 2009 prepublication text, and October 2009 JEBO final are distinct records. The earlier two records were not read in full for this comparison; the final publisher preview was inspected, but its authenticated complete text remains unavailable.

In the inspected model, headquarters allocates a fixed capital total between an established division and an uncertain division. An interested division can pay to turn a favorable private recommendation about the other division into an unfavorable one; a separate noisy public signal cannot be distorted. Headquarters commits ex ante to a signal-contingent investment schedule. Equations (6) and (7) define expected firm values under deterrence and allowance of influence; (8)--(10) optimize the deterrence schedule subject to an influence incentive constraint and fixed total capital. The allowance problem reverses that incentive constraint. Appendix A solves binding-constraint allocation regions and compares the resulting values and signal sensitivities. Rational influence, costly distortion, second-best contracts and the comparison of deterrence with allowance are therefore inherited. Wulf's allowance analysis also uses conditional expected returns: absence of Bayesian reasoning is not a distinction.

Our contrast is the commitment object and hidden action technology. We commit only to factual checking and the stated metric menu; funding subsequently maximizes conditional expected value. The sender can privately purchase either falsification or accurate selection, and a passed packet pools these with genuine good records. Wulf's public signal supplies additional undistorted information about value; replacing its precision by an audit probability does not create our privately costly accurate-selection option or its endogenous submission costs. This comparison establishes a specific difference between inspected mechanisms, not absence of overlap in an unread final or priority over the literature.

Stein (2002), whose author-hosted final article was retrieved, studies the interaction of hard versus soft information with research incentives and capital authority. In particular, its hard-information discussion recognizes productive efforts to influence headquarters. Our quality distribution is fixed and initial information production is a costed required process, so the paper does not reproduce or improve Stein's endogenous productive-research mechanism. Our distinction is between accuracy of a measured object and control of the measured object.

Milgrom and Roberts (1986) is the foundational interested-party disclosure comparison. It studies verifiable communication, sophisticated interpretation and conditions for useful revelation. We use a constrained feasible-evidence technology, with costly factual falsification as an additional option and costly selection under an unobserved recipe. Verifiability alone is therefore not completeness. A scan of the final article was obtained from the author's Stanford archive; the 1985 Cowles record is retained separately in the source log.

Stein (1989) is a directly inspected signal-jamming predecessor. Its Section II explicitly lets the market anticipate equilibrium manipulation while managers still have incentives to manipulate. Consequently, the claim that inference removes all inefficient influence is false in general, and the persistence of influence under rational inference is not our novelty. Here the ambiguity is a common positive packet generated by distinct evidence actions, rather than the dynamic earnings process in that article.

Recent evidence-selection work makes the novelty bar higher. Farina et al. (2026) studies selection versus concealment in verifiable disclosure. The five-author October 16, 2025 manuscript was read for its model, maximally selective equilibria and policy discussion; the final publisher record identifies July 10, 2026 publication and DOI rdag076. We do not treat the 2024 four-author draft as the final paper. That manuscript has a sender observing a state and $N$ signals, disclosing up to $K$ verifiable signals to a receiver who chooses an action. Its maximally selective equilibrium and concealment comparison make selection under rational inference prior art. Fabrication is infeasible there; our narrow distinction is a costly falsification instrument that factual verification removes while leaving costly selection. The present boundary theorem limits, rather than generalizes, the policy reversal. Their experimental findings are theirs; no such experiment is reported here.

Boleslavsky, Carlin, and Cotton (2017), \emph{Competing for Capital}, was now retrieved as the complete March 2017 NBER Working Paper 23273. Its Section 2 has two agents choosing publicly observable, noncontractible quality, then an observable reporting/monitoring policy before privately learning project value. The allocator uses Bayesian posteriors and allocates sequentially without transfers. Proposition 1 characterizes full revelation versus mixed exaggeration; Proposition 3 shows why prohibiting exaggeration can reduce the principal's payoff by discouraging quality investment. Thus noncommitted Bayesian funding and evaluation-induced incentives are already present, not distinctive features of our paper.

Their distinct February 26, 2019 manuscript, \emph{Disincentive Effects of Evaluation}, was retrieved from the current author page. Section 3 has a continuum of agents with observed quality investment, heterogeneous quadratic costs and binary evaluations with no false negatives for high value. Evaluation detects low value with probability $\gamma_i$, and funding capacity is a fixed mass. Propositions 2--4 compare full, absent and agent-designed evaluations; the loss from more informative evaluation operates through lower productive investment. We neither identify the two manuscripts as one final article nor infer final publication from a working-paper record.

Our proposed increment relative to these inspected mechanisms is narrower: after fixed quality is learned, the sender privately chooses between costly falsification and costly factually valid selection; headquarters commits to random checking while the allocator uses the stated pooled signal partition. The new results concern the exact displacement and resource frontier, the exclusion of mixed-regime interior maxima, and finite global intensity candidates. Randomized monitoring as a fraud deterrent, rational influence, second-best contracts and incentive effects of evaluation are inherited principles. A new example does not establish broad priority. The inspected Wulf prepublication mechanism permits the bounded comparison above; the unavailable authenticated final cannot establish absence of an overlapping result.

We read Ba\'nka and Chudziak's \emph{More Claims, Less Evidence} in the primary HTML of arXiv 2610.05547v1 (submitted October 4, 2026). Section 3 fixes the number of false claims in an unreliable artifact and samples uniformly without replacement. In translated notation, with $N_c$ claims, $F_c\geq1$ false claims and $B_c$ checks, its bad-artifact pass probability is
\[
q_c=\frac{\binom{N_c-F_c}{B_c}}{\binom{N_c}{B_c}},
\]
with zero numerator when too few true claims exist. Its Proposition 1 shows that adding true claims with false count and audit budget fixed raises this probability and lowers positive predictive value. Complete checking detects an artifact with at least one false claim. This is finite factual coverage, not the truthful-selection action $M$. Its FEVEROUS/FEVER benchmark claim-set exercises use supplied labels and an oracle verifier; they are neither our computations nor firm evidence. The v1 record's conference-acceptance comment does not authenticate a conference final.

\paragraph{Result-level contribution map.}
Theorem \ref{thm:random_eq}'s relative payoffs (\ref{eq:randompay}) require the available $I/M$ choices; the critical strict-inflation cutoff is $f+\tau rD$, rather than $rD$. Uniform private selection costs then yield the all-submit relation (\ref{eq:randomstrict}). Its $\rho\tau^2D$ term combines selection displaced by auditing with survival of remaining inflation. Passed and submitted masses in (\ref{eq:randommass}) are different endogenous objects, and the real-cost objective (\ref{eq:randomvalue}) charges failed submissions and completed checks too. A public signal or an exogenous false-claim survival probability supplies none of these choice and cost identities by itself.

The mixing frontier and its resource term in (\ref{eq:mixfrontier})--(\ref{eq:mixW}), together with the positive second derivative proved in Appendix \ref{app:random_opt}, use these choices and the uniform private-cost integral. They exclude an interior mixing-regime welfare maximum, while Theorem \ref{thm:random_opt} permits an all-submit interior maximum. Bayesian conditioning, incentive comparisons, probabilistic fraud deterrence, clipping, compactness and stationary-point enumeration are inherited reasoning tools. The theorem's particular polynomial candidates come from the derived frontier, not from a new optimization method. Neither Wulf's public-signal experiment nor the fixed-composition claim-sampling experiment alone supplies those equations; adding a private truthful-selection technology and endogenous submission behavior would be additional modeling work. No inspected source has been shown here to imply the complete frontier, but this targeted comparison is not a broad absence or first-in-literature certificate. The distinct Boleslavsky--Carlin--Cotton comparisons above remain essential: sequential Bayesian allocation and evaluation-induced incentives are already theirs.

The nearest retrieved experiment on structured capital-budgeting aids remains Hartmann and Wei\ss enberger (2024, online 2023). It compares rules and principles, not factual verification crossed with control of the metric universe. Nikias, Schwartz, and Young (2021) concerns information about competing project costs. The primary search was targeted, not a systematic census. No top-journal readiness or secure substantial originality follows from the present comparison.

\section{Model: private quality, private tailoring cost and feasible evidence}\label{sec:model}
\subsection{Capital and values}
One candidate requires one indivisible unit of capital. Headquarters can fund it or retain the capital for a fallback normalized to zero. Project value net of implementation cost and the baseline capital opportunity cost is $v>0$ when $\theta=G$ and $-\ell<0$ when $\theta=B$. The prior good probability is $p\in(0,1)$. Values are risk-adjusted expected net surplus primitives; realized project payoffs need not be deterministic in an application. No return residual defines quality.

An active evaluation architecture first generates a mandatory baseline record at resource cost $\kappa\geq0$ per registered opportunity. The process informs the manager about $\theta$. A good project automatically generates a positive canonical record. A bad project automatically generates a negative canonical record. The process also gives a bad manager private knowledge of the resource cost $c$ of creating a factually accurate favorable selection. Conditional on $B$, $c$ has common-knowledge CDF $F$ on $[0,\bar c]$, continuous and strictly increasing, with positive density in the interior. The allocator observes neither $\theta$ nor $c$ nor the hidden selection recipe.

The mandatory-report assumption fixes the evidence support: a good project produces a positive report without a discretionary application decision. Good managers have no free action that suppresses or relabels this record. This can represent a signed internal register with automatically transmitted positive results. It is a substantive restriction, not an equilibrium refinement. Allowing costless voluntary suppression would create possible non-disclosure equilibria and is outside our uniqueness claim. Bad managers can preserve the negative default or advocate through either action below.
\subsection{Two evidence actions and what the audit sees}\label{subsec:technology}
A bad manager can choose:
\begin{enumerate}
\item $N$: transmit the unchanged negative record, with no additional preparation cost;
\item $I$: inflate the canonical factual value into a positive packet, at real cost $f>0$;
\item $M$: assemble an accurate positive statistic from a favorable selection, at real private cost $c$.
\end{enumerate}
Both $I$ and $M$ produce the same positive packet format. The relevant measurement universe, omitted comparison group or attribution recipe is not observable to the allocator under open metrics. The $M$ packet is factually true about its selected object, but that object does not establish whole-project value. A numerical example is a favorable benefit calculation on a selected subset whose full eligible population is not in the packet. The model deliberately gives the allocator no additional recipe identifier that distinguishes this from an ordinary favorable canonical result. If such provenance were freely available, this information structure would not apply.

A factual verifier checks whether a submitted assertion is true about its submitted object. Complete checking catches $I$ with probability one and does not reject $M$. The verifier is not corrupt or fooled by arithmetic; its mandate and observations do not certify omitted-universe relevance. Its public output is pass or fail. Thus ``full audit'' means a complete factual check of the packet, not a complete audit of economic relevance. A good packet always passes.

A metric-governance protocol locks an enforceable canonical object and its population definition before quality and tailoring cost are observed. Its effect is to remove $M$ from the feasible action set. It does not detect factual inflation without the audit. The protocol costs $K\geq0$ per registered opportunity. Enforceability is assumed and must be established in an application; the model does not derive a cheap way to police every economically complex metric. This separation of factual and recipe observability is the key institutional assumption replacing the previous fixed score.
\paragraph{Observation: factual coverage and relevance under the maintained partition.}
Fix a submitted packet $x$ whose claims are all true. Suppose two possible hidden universe/recipe completions differ in economic relevance but produce the same packet and verifier-accessible information. If the checker chooses claims using only that information, previous truth answers and independent randomness, then increasing coverage, even to complete checking, cannot distinguish those completions. Conditional on $x$, its transcript supplies no additional relevance information under this partition.

\emph{Proof.} Couple the checks under the two completions using the same random seed. The first query is identical and receives the same true answer. Induction gives identical subsequent queries and answers, including under adaptive stopping. Thus the conditional transcript distributions coincide; Bayes' rule leaves the conditional relevance odds unchanged wherever defined. This does not say that $x$ itself is uninformative. Observing the omitted universe or recipe can refine the information partition and invalidate the equivalence. This is a short implication of our maintained technology, not a general impossibility theorem or the paper's frontier contribution. In particular, the source's unreliable-artifact case $F_c\geq1$ is different: complete factual coverage detects its false claim, whereas $M$ has no false submitted claim.\hfill$\square$

\subsection{Payoffs, timing and menu}\label{subsec:timing}
The manager receives private control benefit $r>0$ if funded, less the chosen preparation resource cost. Benefits $r$ are not added to the firm's objective. Costs $f,c$ represent the manager's opportunity time and are also losses of real firm resources in the same numeraire. They are counted once in firm net value. We are not computing a utilitarian welfare sum that separately values managerial status.

Headquarters chooses an architecture $(A,L)\in\{0,1\}^2$: factual auditing is absent or complete, and metric choice is open or canonically locked. In the zero-shock baseline it can instead decline to generate evidence and choose the fallback, obtaining zero. With the outside-option shock, declining evidence retains the optimal prior-based funding decision characterized in Section \ref{sec:global}. This is the binary policy menu used in Sections \ref{sec:eq}--\ref{sec:global}. Section \ref{sec:random} enlarges open-metric auditing to a committed continuous intensity, leaving the locked menu unchanged. No claim of unrestricted mechanism optimality is made.

After policy choice and baseline production, the manager learns $\theta,c$ and selects evidence. The verifier acts if $A=1$. The allocator sees a negative record $n$, a passed positive packet $+$, or a failed audit $d$, and then chooses funding. It does not commit to a funding probability. It uses posterior expected project value, taking past preparation and verification costs as sunk. The policy stage anticipates those costs and equilibrium responses.

Incremental packet evaluation costs $e\geq0$ and full factual checking costs $a\geq0$ are incurred per positive packet submitted. Routine handling of every negative baseline record is included in $\kappa$. Failed positive packets still incur both packet costs. Initial information production, preparation, evaluation, audit and governance are distinct resource accounts. As every realized good record is positive, the positive-packet event always has positive probability.

\section{Bayesian equilibrium}\label{sec:eq}
Let $u$ be the allocator's probability of funding after $+$. For a bad type, the three expected payoffs are
\begin{equation}U_N=0,\qquad U_I=(1-A)ru-f,\qquad U_M=ru-c\quad(L=0).\label{eq:payoffs}\end{equation}
In this zero-shock section the audit failure branch is unfunded; Section \ref{sec:global} allows funding after a detected failure when the fallback is worse. Denote by $s$ the conditional probability that a bad project produces a passed positive packet. Bayes' rule and the funding rule are
\begin{equation}\mu(s)=\Pr(G\mid+)=\frac{p}{p+(1-p)s},\qquad
u(s)=\frac{pv-(1-p)\ell s}{p+(1-p)s}.\label{eq:posterior}\end{equation}
The allocator funds if $\nu(s)>0$, rejects if $\nu(s)<0$, and may randomize only if $\nu(s)=0$. Write
\begin{equation}T=\frac{pv}{(1-p)\ell},\quad \alpha=F(f),\quad c_T=F^{-1}(T),\quad P=p+(1-p)T.\label{eq:primitives}\end{equation}
The baseline theorem concerns the nonempty open region
\begin{equation}0<f<r,\qquad F(f)<T<F(r)<1.\tag{R}\label{eq:region}\end{equation}
In particular, the prior expected project value is negative, $0<T<1$, and $f<c_T<r<\bar c$. Continuous parameters need not be tuned to an equality. The posterior indifference below is an endogenous mixed-equilibrium condition.
\begin{theorem}[All equilibria in region R]\label{thm:eq}
Within the mandatory-record game, every perfect Bayesian equilibrium has the following funding probabilities and aggregate bad-manager actions:
\begin{center}\small
\begin{tabular}{lrrrr}\toprule
Policy & $u$ & Metric $M$ & Inflation $I$ & Negative $N$\\\midrule
Open, no audit $(0,0)$ & $f/r$ & $\alpha$ & $T-\alpha$ & $1-T$\\
Open, full audit $(1,0)$ & $c_T/r$ & $T$ & $0$ & $1-T$\\
Locked, no audit $(0,1)$ & $f/r$ & $0$ & $T$ & $1-T$\\
Locked, full audit $(1,1)$ & $1$ & $0$ & $0$ & $1$\\\bottomrule
\end{tabular}
\end{center}
Action masses are conditional on a bad project. For $(0,0)$, all $c<f$ choose $M$, and types $c>f$ distribute between $I$ and $N$ with aggregate inflation mass $T-\alpha$. One implementation inflates with constant probability $(T-\alpha)/(1-\alpha)$ among $c>f$. For $(1,0)$, types $c<c_T$ choose $M$, and types $c>c_T$ choose $N$. With $(0,1)$, any type allocation of inflation with total mass $T$ works.

For the first three policies, $s=T$, $\mu=\ell/(v+\ell)$ and $\nu=0$. Negative records have posterior good probability zero. Under $(1,1)$, positive and negative records reveal quality. Failed factual audits have good probability zero by technology, including off equilibrium path. There is multiplicity in which indifferent bad types inflate without an audit, but not in aggregate costs, posteriors, funding rates or firm value. Types exactly at a cutoff have zero probability. No additional outcome equilibrium exists under the specified support.
\end{theorem}
The proof in Appendix A examines every receiver probability and sender deviation. Thus the displayed actions are not fitted to an imposed funding rule. The allocator's posterior explicitly discounts strategic reporting, yet a positive packet cannot reveal which private-cost manager selected it.

Full factual checking produces no detected falsifications in equilibrium: managers anticipate detection and switch instruments or preserve the negative default. Consequently, a zero detection rate does not establish that the project packet is economically representative. Nor does it establish that auditing has no incentive effect. Here that effect is a complete shift in the composition of influence.
\begin{corollary}[Residual allocation loss after inference]\label{cor:regret}
For each of the first three policies, expected project surplus before information costs is zero. Full-information surplus is $pv>0$. Their allocation regret is therefore $pv$, including both funding bad projects and rejecting good ones. Full audit under open metrics changes the good and bad funding probabilities from $(pf/r,(1-p)Tf/r)$ to $(pc_T/r,(1-p)Tc_T/r)$, increasing both. Full audit with locked metrics funds exactly the good project and has zero allocation regret.
\end{corollary}
This is residual uncertainty relative to full information, not biased beliefs. Under the open-audit policy there is no deterministic influence constant to subtract. Subtracting the average prevalence of selection changes the interpretation of the packet but does not identify the state of an individual manager's project. The equilibrium already incorporates that prevalence through (\ref{eq:posterior}).

\section{Costs and optimal governance}\label{sec:policy}
Define the private-cost truncated integral and two nonnegative quantities
\begin{equation}I(x)=\int_0^x c\,dF(c),\quad J=I(f)+f[T-F(f)],\quad D=fF(f)-I(f)=\int_0^f(f-c)\,dF(c).\label{eq:costs}\end{equation}
The word ``private'' describes who knows the type, not an extra cost category. $(1-p)I(x)$ or $(1-p)J$ already counts the full real preparation cost for a registered opportunity.
\begin{theorem}[Exact substitution cost and complementarity]\label{thm:policy}
In region R the four active policies have firm net values
\begin{align}
W_{00}&=-\kappa-eP-(1-p)J,\label{eq:w00}\\
W_{10}&=-\kappa-(e+a)P-(1-p)I(c_T),\label{eq:w10}\\
W_{01}&=-\kappa-eP-(1-p)fT-K,\label{eq:w01}\\
W_{11}&=pv-\kappa-(e+a)p-K.\label{eq:w11}
\end{align}
With open metrics, the effect of full auditing is
\begin{equation}W_{10}-W_{00}=-(1-p)\int_f^{c_T}(c-f)\,dF(c)-aP<0.\label{eq:auditloss}\end{equation}
With no audit, locking the metric has effect
\begin{equation}W_{01}-W_{00}=-(1-p)D-K<0.\label{eq:lockloss}\end{equation}
The finite-menu interaction is strictly positive:
\begin{equation}\begin{split}
\Gamma&=W_{11}-W_{01}-W_{10}+W_{00}\\
&=pv+(e+a)(1-p)T+(1-p)[I(c_T)+D]>0.
\end{split}\label{eq:interaction}\end{equation}
Among active architectures, the optimum is either $(0,0)$ or $(1,1)$. The package wins exactly when
\begin{equation}K<K^*_{\rm active}=pv+e(1-p)T-ap+(1-p)J.\label{eq:active}\end{equation}
Including the feasible no-evidence fallback with value zero, the optimal policy is $(1,1)$ exactly when
\begin{equation}p(v-e-a)>\kappa+K,\label{eq:optimal}\end{equation}
and is the fallback when the inequality is reversed. Equality permits both. The other active policies have strictly negative value and cannot beat the fallback.
\end{theorem}
Equation (\ref{eq:auditloss}) says more than that audit expenditure is costly. Even at $a=0$, full factual checking worsens firm value through higher real preparation costs. The marginal types shifted to selection have costs above $f$ and previously either inflated at cost $f$ or preserved the negative record. Their new costly selection is a private best response to the receiver's endogenously higher funding probability. The initial-cost increase $\kappa$ and evaluation cost $e$ are not counted again inside the tail integral.

Equation (\ref{eq:interaction}) is an intervention complementarity, not a claim of globally optimal governance. The information-production technology, admissible evidence actions and menu are restricted. Arbitrary outcome-contingent transfers, penalties for false reports, revelation of selection provenance or fractional audit schedules could alter the solution. In particular, an unrestricted commitment to deny funding after every positive packet could save some later costs but would be a different mechanism. The allocator in this game does not possess such funding commitment.

The fallback comparison is essential. Ignoring it would recommend an active evidence architecture even when every available architecture destroys value. Mandatory reporting is conditional on choosing an active architecture; it does not force the firm to purchase initial information when the option has negative net value.

\section{The local strict-receiver result retained from R1}\label{sec:strict}
The baseline has posterior indifference on a positive packet, as in many binary persuasion games. To show that the reversal is not generated solely by an exogenous tie-breaking rule, let the fallback have privately observed value $t\sim U[-\delta,\delta]$, independent of quality and tailoring cost, realized for the allocator after reporting. Relative to always using this fallback, candidate funding produces value $v-t$ or $-\ell-t$. The manager does not observe $t$ before choosing evidence. The allocator strictly funds if $\nu(s)>t$, apart from a probability-zero equality.

For $|\nu(s)|<\delta$, the induced funding probability is $u=[\nu(s)+\delta]/(2\delta)$. Define
\begin{equation}H_\delta(u)=\frac{p[v-\delta(2u-1)]}{(1-p)[\ell+\delta(2u-1)]}.\label{eq:hdelta}\end{equation}
\begin{proposition}[Solved strict-decision perturbation]\label{prop:strict}
For every primitive vector in region R, sufficiently small $\delta>0$ gives the following unique aggregate equilibrium outcomes. Let $u_0=f/r$ and $s_0=H_\delta(u_0)$. Without an audit, open metrics have metric mass $F(f)$, inflation mass $s_0-F(f)$ and negative mass $1-s_0$; locked metrics have inflation mass $s_0$. With an audit and open metrics, the unique $u_A\in(u_0,1)$ solves
\begin{equation}F(ru_A)=H_\delta(u_A),\label{eq:strictroot}\end{equation}
and bad managers select metrics iff $c<ru_A$. With an audit and locked metrics, only good projects are funded. Both without-audit outcomes retain type-level inflation multiplicity with fixed aggregate mass.

For an imperfect-policy equilibrium $(u,s)$, gross expected value relative to the fallback is
\begin{equation}G_\delta(u,s)=[p+(1-p)s]\delta u^2.\label{eq:gdelta}\end{equation}
Subtract the same five resource categories as before, using the new packet probability and new preparation cutoffs. The solutions and values converge to Theorems \ref{thm:eq}--\ref{thm:policy} as $\delta\downarrow0$. Consequently audit loss with open metrics, metric-locking loss without audit, strict complementarity and the optimal fallback-or-package conclusion all persist for sufficiently small $\delta$. The package value and its threshold (\ref{eq:optimal}) are unchanged.
\end{proposition}
Equation (\ref{eq:strictroot}) is a solved monotone scalar equation for arbitrary $F$; the left side increases and the right side decreases. It uniquely determines actions and beliefs. Thus this extension does not keep the old manager strategy fixed. Numerical examples below solve the root anew and integrate the receiver's strict payoff directly.

The perturbation does not make the entire model unrestricted. Managerial inflation still mixes among cost-indifferent types when it is used. The mandatory good-report support and hidden-recipe technology remain. It does establish that a rational allocator need not be forced to toss a coin at a common tie in every realized evaluation for the mechanism to survive.

\section{The entire uniform-cost outside-option game}\label{sec:global}
The main extension now permits every finite $\delta>0$, not just a neighborhood of zero. Maintain region R and specialize $F(c)=c/C$ on $[0,C]$. Write $\rho=r/C$, $b=f/r$, $\alpha=f/C=\rho b$, and $m=pv-(1-p)\ell<0$. Thus $0<\alpha<T<\rho<1$. The fallback shock is independent of all project variables. Negative incremental fallback values are admissible: adding a common positive level to both opportunities changes no decision or welfare difference. The model imposes no upper bound on $\delta$.

Define the funding and option-value functions, for $\delta>0$,
\begin{align}
h_\delta(x)&=\max\{0,\min\{1,(x+\delta)/(2\delta)\}\},\label{eq:clip}\\
B_\delta(x)&=\begin{cases}0&x\leq-\delta,\\(x+\delta)^2/(4\delta)&-\delta<x<\delta,\\x&x\geq\delta.\end{cases}\label{eq:option}
\end{align}
The correctly optimized no-evidence value, relative to always using the fallback, is
\begin{equation}V_\varnothing=B_\delta(m)=\begin{cases}0&\delta\leq-m,\\(m+\delta)^2/(4\delta)&\delta>-m.\end{cases}\label{eq:none}
\end{equation}
Without evidence the allocator still observes $t$ and funds precisely when $m>t$. No initial or packet costs are incurred. In particular, a wide outside-option distribution does not imply compulsory shutdown.

At a negative record or a detected falsification, the posterior good probability is zero, and funding occurs with probability $n=h_\delta(-\ell)$. This is positive when $\delta>\ell$: a known bad project can dominate an even worse fallback. At a passed positive record, Bayes' rule remains (\ref{eq:posterior}); write
\begin{equation}D(s)=h_\delta(\nu(s))-n.\label{eq:increment}
\end{equation}
This continuous nonincreasing function is the incremental funding benefit of a passed packet. Bad-manager payoffs relative to preserving the negative default are $0$, $rD(s)-f$ for unaudited inflation, $-f$ for audited inflation, and $rD(s)-c$ for metric selection. There is no fixed political premium.

\begin{theorem}[Full-range equilibrium correspondence]\label{thm:global}
For every $\delta>0$, audited open metrics have a unique equilibrium passed-bad mass $y=\rho d$ and metric mass $y$, with zero inflation, where
\begin{align}
A_2&=2\delta(1-p)\rho,\quad
A_1=2\delta p+(\ell-\delta)_+(1-p)\rho,\notag\\
A_0&=p[v+\min\{\delta,\ell\}],\qquad
d=\frac{2A_0}{A_1+\sqrt{A_1^2+4A_2A_0}}.\label{eq:quadratic}
\end{align}
Equivalently, $y=\rho D(y)$. Metric selection occurs exactly for $c<Cy$. Audited locked metrics have $s=M=I=0$ and reveal quality. The receiver nevertheless funds bad quality with probability $n$ and good quality with $h_\delta(v)$.

For an unaudited architecture let $z=\alpha$ under open metrics and $z=0$ under locked metrics. The following cases are exhaustive, with equalities assigned to the middle row:
\begin{center}\small
\begin{tabular}{p{.25\textwidth}p{.25\textwidth}p{.37\textwidth}}\toprule
Condition & Passed-bad mass & Bad-manager actions\\\midrule
$D(1)>b$ & $s=1$ & $M=z$, $I=1-z$, $N=0$\\
$D(z)\geq b\geq D(1)$ & every $s\in[z,1]$ with $D(s)=b$ & $M=z$, $I=s-z$, $N=1-s$\\
$D(z)<b$ & $s=y$ if open; $s=0$ if locked & $M=s$, $I=0$, $N=1-s$\\\bottomrule
\end{tabular}\end{center}
In the third open case, $y<\alpha$. In the middle row, if $n+b<1$, its unique solution is
\begin{equation}
x=\delta[2(n+b)-1],\qquad s=\frac{p(v-x)}{(1-p)(\ell+x)}.\label{eq:inverse}
\end{equation}
If $n+b=1$, all $s$ in
\begin{equation}[z,\min\{1,p(v-\delta)/((1-p)(\ell+\delta))\}]\label{eq:plateau}
\end{equation}
are equilibria whenever that interval is nonempty. This upper-clipping multiplicity can occur with locked metrics in R, but cannot occur with open metrics in R. If $n+b>1$, the middle row is impossible. Types below $f$ select metrics in the middle and pooling open cases. Remaining types distribute between the displayed inflation and negative masses whenever indifferent. There are no other feasible-message equilibria.
\end{theorem}
This is a piecewise solution over the entire shock range: inequalities of explicit clipped functions select the case, (\ref{eq:inverse}) gives the mixing mass, and (\ref{eq:quadratic}) gives the only nonlinear root in closed form. It also covers transition equalities, vanishing inflation and fully positive pooling. Negative records can have zero on-path probability in the pooling case; their off-path posterior is still structurally bad because good types cannot produce them. Failed audits have the same posterior, including deviations. A passed packet always has positive probability because $p>0$.

For orientation, with $p=.4$, $v=\ell=2$, $r=1$, $f=.2$, $C=1.2$, unaudited open metrics are mixed for $0<\delta<2/3$, pooled for $2/3<\delta<4$, mixed again for $4<\delta<8$, and metric-only for $\delta>8$. Unaudited locked metrics have the same first three regimes, with the second mixing regime extending to $\delta=10$, after which there is no advocacy. Endpoint outcomes join continuously. At $\delta=2$ negative and failed records begin to receive positive funding probability for larger shocks. At $\delta=.4$ the no-evidence value begins to be positive. These are distinct boundaries, not interchangeable approximations.

\subsection{All information sets, costs and governance}\label{subsec:globalcost}
Let $M,I$ denote bad-type action masses, also for a counterfactual deviation profile. Passed positives have bad mass $s=M+(1-A)I$; failed positives have unconditional probability $q_d=(1-p)AI$; unchanged negatives have probability $q_n=(1-p)(1-M-I)$. Submitted positive packets have probability $Q=p+(1-p)(M+I)$, whereas passed positives have probability $P(s)=p+(1-p)s$. The complete accounting identity is
\begin{align}
G&=P(s)B_\delta(\nu(s))+(q_n+q_d)B_\delta(-\ell),\label{eq:allsets}\\
W_{AL}&=G-\kappa-(e+Aa)Q-(1-p)[CM^2/2+fI]-LK.\label{eq:fullcost}
\end{align}
The selection integral assumes the equilibrium lowest-cost interval, as in Theorem \ref{thm:global}; the formula also applies to deviations preserving that interval. For arbitrary individual deviations replace $CM^2/2$ by the actual integral over selecting types. Failed packets incur preparation, evaluation and audit costs even if funded after failure. On path, audited inflation is strictly dominated, so $I=q_d=0$; it is not legitimate to omit these costs from the deviation calculation. Negative-record processing is in $\kappa$ once. $r$ is a private incentive and is not added to firm surplus.

In equilibrium define
\begin{equation}\mathcal G(s)=P(s)B_\delta(\nu(s))+(1-p)(1-s)B_\delta(-\ell).\label{eq:gfull}
\end{equation}
Then (\ref{eq:fullcost}) with $G=\mathcal G(s)$ evaluates every architecture. Full-information value is $\mathcal G(0)=pB_\delta(v)+(1-p)B_\delta(-\ell)$, not necessarily $pv$. Its difference from $\mathcal G(s)$ is the gross information regret, nonnegative by Blackwell refinement. For finite $\delta>0$, a metric-only equilibrium has $y>0$; its passed packet does not reveal individual quality despite rational inference.

\begin{proposition}[Optimal governance, including genuine multiplicity]\label{prop:menu}
Let $O_0=W_{00}$ and $O_1=W_{10}$ from (\ref{eq:fullcost}). Let $L_0=W_{01}+K$ and $L_1=W_{11}+K$ be the locked values before governance cost. For any specified equilibrium in the unaudited locked correspondence, the optimal value is
\begin{equation}\max\{V_\varnothing,O_0,O_1,L_0-K,L_1-K\}.\label{eq:menu}
\end{equation}
An optimal locked policy exists precisely when
\begin{equation}K\leq\max\{L_0,L_1\}-\max\{V_\varnothing,O_0,O_1\}.\label{eq:kglobal}
\end{equation}
Within locked policies, choose full audit when $L_1\geq L_0$. Within open policies choose full audit exactly when
\begin{equation}
a\leq a^*:=\frac{\mathcal G(y)-\mathcal G(s_0)+C_0-C_A+e[P(s_0)-P(y)]}{P(y)},\label{eq:afees}
\end{equation}
where $s_0$ is the unaudited open mass, $C_0$ its unconditional preparation cost, and $C_A=(1-p)Cy^2/2$. Ties permit both choices. On the locked clipping interval (\ref{eq:plateau}), $W_{01}$ is affine in $s$, so its two endpoints give the exact possible-value interval. A selection-free policy claim must hold against the relevant endpoint, rather than silently choosing a favorable equilibrium.
\end{proposition}
These are evaluable thresholds, not an optimization over unsolved strategies. Every input is a closed-form function of primitives and the stated piecewise case. The menu remains discrete: transfers, fines, intermediate audit intensity and freely observed provenance are not feasible policies here.

\subsection{A boundary on the policy reversal}\label{subsec:boundary}
\begin{theorem}[No costless-audit reversal from a valuable open architecture]\label{thm:boundary}
For uniform costs in R, any finite $\delta\geq0$ and nonnegative resource costs, suppose the unaudited open architecture satisfies $W_{00}\geq V_\varnothing$. Holding all other costs fixed, full factual audit at $a=0$ weakly improves its value. Improvement is strict when unaudited inflation has positive mass. Consequently, no nonempty region has both a valuable unaudited open architecture and a negative net effect of costless factual verification caused by metric substitution alone. The threshold $a^*$ in (\ref{eq:afees}) is nonnegative under that hypothesis, and is positive with inflation.

There do exist open sets with a strictly optimal imperfect active architecture. There also exist open sets with a strictly optimal unaudited imperfect architecture, increased metric preparation after audit, and a negative net audit effect at positive $a$. In the latter case audit expenditure is necessary for the negative net effect; substitution alone cannot explain it.
\end{theorem}
This narrows the contribution. The baseline resource-dissipation reversal is a correct conditional comparison of architectures that are globally dominated there. Large shocks break the perfect-versus-no-evidence dichotomy, but do not rescue that strong reversal at a valuable starting architecture. Calling the positive-fee example a pure substitution welfare loss would misstate the result.

\subsection{Synthetic witnesses with strict governance margins}\label{subsec:witnesses}
All following values are synthetic, in arbitrary units. Example A uses $p=.4$, $v=\ell=2$, $r=1$, $f=.2$, $C=1.2$, $\kappa=.02$, $e=.03$, $a=.04$, $K=1$, $\delta=1$. Example B changes $r=.02$, $f=.004$, $C=.024$, $\kappa=e=.0001$, $a=.1$, $K=1$, $\delta=.1$, retaining $p,v,\ell$. Both satisfy R with positive administrative costs.
\begin{center}\small
\begin{tabular}{lrr}\toprule
Optimized option & Example A & Example B\\\midrule
No evidence & .090000 & .000000\\
Open, no audit & -.070000 & .001618\\
Open, full audit & .109885 & -.036056\\
Locked, no audit & -1.080000 & -.998582\\
Locked, full audit & -.248000 & -.240140\\\bottomrule
\end{tabular}\end{center}
In A, audited open metrics are uniquely optimal within the binary-audit menu. The expanded menu in Section \ref{sec:random} improves on full auditing. The equilibrium has $y=.519812$, $u=.623774$, posterior good probability $.561887$, gross value $.276991$ and preparation cost $.097273$. This is a valuable imperfect architecture, with a strict gain $.019885$ over optimally using no evidence.

In B, unaudited open metrics are uniquely optimal within the binary menu; Section \ref{sec:random} also checks the enlarged intensity menu. Audit raises metric mass from $.166667$ to $.632939$ and preparation cost from $.001499$ to $.002884$. It also raises gross value from $.003299$ to $.044983$. The costless-audit net gain is $.040303$, including saved evaluation costs. The break-even audit fee is $a^*=.051686$ per submitted packet. The actual fee $.1$ reverses the net comparison. This establishes coexistence of valuable imperfect governance and costly metric substitution, but the loss requires verification expenditure. Strict margins and continuity yield neighborhoods of both witnesses, not isolated equalities.

\section{Endogenous random factual auditing under open metrics}\label{sec:random}
The earlier theorems retain their stated binary-audit menu. We now enlarge only its open-metric option: headquarters commits before information production to an independent check of each submitted positive packet with probability $\tau\in[0,1]$. Each completed check costs $a$. A checked inflation fails, while unchecked inflation and accurate selected packets retain the positive signal. There are no transfers or penalties, no new data and no productive-quality investment.

To implement the stipulated common positive posterior, the allocator observes only unchanged negative, failed, or surviving positive. Checked-pass and unchecked-pass flags are pooled before the allocation decision. Detailed scheduling and billing records are not available to the allocator at that point. Headquarters commits to the checking process, not to funding. If checked-pass flags were separately observable, there would be two positive posteriors and a different game; none of the single-$D$ formulas below should be applied to that information structure. This restriction makes the observation technology explicit.

Continue with uniform costs in R. Use $\rho=r/C$, $b=f/r$, $\alpha=f/C=\rho b$, $n=h_\delta(-\ell)$ and $D(s)=h_\delta(\nu(s))-n$ when $\delta>0$. A checked failure is structurally bad and can still be funded with probability $n$. Bad-manager payoffs relative to the negative default are
\begin{equation}
0,\qquad (1-\tau)rD(s)-f,\qquad rD(s)-c.\label{eq:randompay}
\end{equation}
Let $M,I,N$ be conditional bad action masses. Passed bad mass and all submitted-positive probability differ:
\begin{equation}s=M+(1-\tau)I,\qquad Q=p+(1-p)(M+I).\label{eq:randommass}\end{equation}
Bayes' rule uses $s$, while packet and expected audit expenditure use $Q$. Full accounting, including completed failed checks, is
\begin{equation}
W(\tau)=\mathcal G(s)-\kappa-eQ-a\tau Q-(1-p)[CM^2/2+fI].\label{eq:randomvalue}
\end{equation}
Here $\mathcal G$ is (\ref{eq:gfull}). Failures have probability $(1-p)\tau I$; together with unchanged negatives they contribute $(1-p)(1-s)B_\delta(-\ell)$. The private benefit $r$ is never added to firm value.

\subsection{All intensity equilibria and the exact deterrence boundary}\label{subsec:random_eq}
\begin{theorem}[Random-audit equilibrium, all positive shock sizes]\label{thm:random_eq}
For every $\delta>0$ and $\tau\in[0,1]$, there is exactly one aggregate open-metric equilibrium. Let $y$ be the metric-only root (\ref{eq:quadratic}), $d_y=y/\rho$, and
\begin{equation}\tau_* =\max\{0,1-b/d_y\}\in[0,1).\label{eq:taustar}\end{equation}
Inflation is absent exactly when $\tau\geq\tau_*$. On this interval $s=M=y$, $I=0$, $N=1-y$ and the positive funding probability is $n+d_y$.

For $\tau<\tau_*$ define $d_0=b/(1-\tau)$, $M_0=\rho d_0$, and $U=1-\tau+\tau M_0$.
\begin{enumerate}
\item If $D(U)\leq d_0$, inflation and the negative default tie. Set
\begin{equation}
x=\delta[2(n+d_0)-1],\quad s=\frac{p(v-x)}{(1-p)(\ell+x)},\quad M=\rho d_0,\quad I=\frac{s-M}{1-\tau},\quad N=1-M-I.\label{eq:randommix}
\end{equation}
Then $0<I\leq1-M$, with equality at $D(U)=d_0$. Types $c<rd_0$ select metrics. Types above that cutoff distribute the displayed inflation and negative masses in any measurable manner.
\item If $D(U)>d_0$, every bad type submits a positive packet. The unique root
\begin{equation}s=1-\tau+\tau\alpha+\rho\tau^2D(s)\label{eq:randomstrict}
\end{equation}
gives $d=D(s)$, $M=\alpha+\rho\tau d$, $I=1-M$, $N=0$. Selection occurs iff $c<f+\tau rd$, not iff $c<rd$. Funding of a passed packet is $n+d$. It can equal one, and that upper clipping is included.
\end{enumerate}
There is no lower-clipped positive funding equilibrium. The mixed case is always strictly inside the receiver response interval. At equality between cases, actions and values agree. There is type-level mixing multiplicity in case 1 but no aggregate or welfare multiplicity, including at $\tau_*$. Negative or failed messages have posterior good probability zero even off path. The positive message has probability at least $p$.
\end{theorem}
The formula for $\tau_*$ is a deterrence threshold, not a welfare optimum. At $\tau_*$ high-cost types may be individually indifferent about a deviation to inflation, but positive aggregate inflation would change the posterior and cannot be an additional equilibrium. For $a>0$, every intensity above $\tau_*$ is strictly worse than $\tau_*$: evidence actions and allocation remain the same while completed audit costs increase. In particular, full factual audit is never optimal within the expanded open-metric option when $a>0$.

\subsection{An explicit finite global candidate set}\label{subsec:candidates}
The full solution allows multiple mixed/all-submit transitions as intensity changes; it does not assume the frontier is monotone. Define constants and functions, for $\delta>0$,
\begin{align}
q&=p/(1-p),&\lambda&=p(v+\ell)/(1-p),&L&=(\ell-\delta)_+,\notag\\
H(d)&=L+2\delta d,&S(d)&=\lambda/H(d)-q,&d_{\max}&=1-n,\label{eq:frontierS}\\
g(d)&=\delta n^2+(1-p)\lambda\delta\,d(d+2n)/H(d).\label{eq:frontierg}
\end{align}
For interior receiver funding, $s=S(d)$ and gross value is $g(d)$. The inverse is used only when $n+d<1$; the upper-clipped case is treated separately. $H(d)>0$ on all feasible interior branches.

In the inflation/negative mixing regime, parameterize by $d$, not $\tau$:
\begin{align}
\tau(d)&=1-b/d,\quad M=\rho d,\quad I=(d/b)[S(d)-\rho d],\notag\\
Q(d)&=(1-p)\{q+\rho d+(d/b)[S(d)-\rho d]\},\label{eq:mixfrontier}\\
W_m(d)&=g(d)-\kappa-[e+a(1-b/d)]Q(d)\notag\\
&\hspace{1cm}-(1-p)r[dS(d)-\rho d^2/2].\label{eq:mixW}
\end{align}
Feasibility requires $b\leq d\leq d_y$ and $\rho d\leq S(d)\leq\rho d+b/d-b\rho$. Thus (\ref{eq:mixW}) is an explicitly restricted rational frontier, not an unconstrained policy objective.

For the all-submit interior regime let
\begin{equation}
E(\tau,d)=[q+1-(1-\alpha)\tau+\rho\tau^2d]H(d)-\lambda=0.\label{eq:frontierE}
\end{equation}
The relevant root is positive and regular. Since $Q=1$, preparation cost simplifies to $(1-p)[f-f^2/(2C)+C\rho^2\tau^2d^2/2]$. With $J=(1-p)C\rho^2$, define
\begin{align}
N_1&=1-\alpha-2\rho\tau d,&
d'&=\frac{N_1}{\rho\tau^2+2\delta\lambda/H(d)^2},\label{eq:dprime}\\
g_d&=\frac{2(1-p)\lambda\delta[L(d+n)+\delta d^2]}{H(d)^2},\notag\\
W_i'(\tau)&=[g_d-J\tau^2d]d'-a-J\tau d^2.\label{eq:strictderivative}
\end{align}
An explicit polynomial equation for its stationary points is
\begin{align}
R(\tau,d)={}&[G_1-J\tau^2dH^2]N_1\notag\\
&-(a+J\tau d^2)[\rho\tau^2H^2+2\delta\lambda]=0,\label{eq:frontierR}\\
G_1={}&2(1-p)\lambda\delta[L(d+n)+\delta d^2].\notag
\end{align}
This polynomial is obtained by multiplying (\ref{eq:strictderivative}) by strictly positive denominators, so no stationary point is lost. Algebraic roots from another branch must still be discarded.

When the receiver funds every surviving positive, $d=d_{\max}$ and $s=1-(1-\alpha)\tau+\rho\tau^2d_{\max}$. On this feasible branch $W$ is a strictly concave quadratic in $\tau$. Its only possible interior maximum is
\begin{equation}
\tau_c=\frac{(1-p)[\ell+B_\delta(-\ell)](1-\alpha)-a}{(1-p)\{2[\ell+B_\delta(-\ell)]\rho d_{\max}+C\rho^2d_{\max}^2\}}.\label{eq:clippeak}
\end{equation}
It is a candidate only if $0<\tau_c<\tau_*$, inflation strictly beats $N$, and $\nu(s)\geq\delta$.

\begin{theorem}[Constructive global intensity optimum]\label{thm:random_opt}
For $\delta>0$, $W(\tau)$ is continuous on $[0,1]$ and attains a global maximum. A complete finite candidate procedure is:
\begin{enumerate}
\item Include $0,\tau_*,1$.
\item Include every feasible mixed/all-submit boundary. Solve the cubic
\begin{equation}\mathcal T(d)=\lambda d-[(q-b\rho)d+\rho d^2+b]H(d)=0,\quad \tau=1-b/d,\label{eq:transitionpoly}
\end{equation}
retaining only actual transition points in $[0,\tau_*]$.
\item Include every feasible clipping boundary from $E(\tau,d_{\max})=0$, a quadratic in $\tau$.
\item Include $\tau_c$ if feasible, and every feasible all-submit interior stationary root of $E=R=0$.
\end{enumerate}
For the last step compute the polynomial resultant $\operatorname{Res}_d(E,R)$, isolate all its real roots in $(0,\tau_*)$, and retain only roots on the actual equilibrium interior branch satisfying (\ref{eq:strictderivative}). If the resultant is identically zero, the interior-branch value is constant and its endpoints suffice. Clearing factors or denominators can create extraneous roots, never additional valid policies. Compare actual $W$ at the retained candidates; this maximum is necessary and sufficient for global optimality.

No strict interior optimum can occur within a nondegenerate inflation/negative mixing interval. An interior optimum with positive inflation can occur in the all-submit regime, including upper clipping. If candidate values tie, retain every maximizer. If $a=0$ and a metric-only candidate maximizes value, every intensity in $[\tau_*,1]$ maximizes it as well; constant interior branches, if present, are retained in full. Other isolated ties require no selection rule.
\end{theorem}
This is polynomial root enumeration across actual regime boundaries, not a numerical grid or a generic optimizer declaration. For rational primitives the reproduction package constructs exact rational coefficients, uses Sturm counts to isolate every real polynomial root, and filters using the equilibrium conditions. Floating arithmetic is used for final root evaluation and displayed values, with tolerances reported. The proofs establish the all-primitive result; the eight executed cases establish implementation checks within those assumed models.

\subsection{The zero-width endpoint}\label{subsec:random_zero}
The original $\delta=0$ game is also solvable for every intensity; it must not be handled by dividing by $\delta$. Set $n=0$, $y=T$, $d_y=T/\rho$, and $\tau_*=1-\alpha/T$. For $\tau\geq\tau_*$, $s=M=T$, $I=0$, with positive funding probability $u=T/\rho$. For $\tau<\tau_*$ let $u_0=b/(1-\tau)$ and $U=1-\tau+\tau\rho u_0$. If $T\leq U$, set $s=T$, $u=u_0$ and use (\ref{eq:randommix})'s action formulas with $d_0=u_0$. If $T>U$, all bad types submit and
\begin{equation}
\widehat u=\frac{T-1+(1-\alpha)\tau}{\rho\tau^2}.\label{eq:zero_u}
\end{equation}
If $\widehat u\leq1$, use $u=\widehat u$, $s=T$; if $\widehat u>1$, use $u=1$, $s=1-(1-\alpha)\tau+\rho\tau^2<T$. In either case $M=\alpha+\rho\tau u$, $I=1-M$. The denominator in (\ref{eq:zero_u}) is used only in this all-submit case, which excludes $\tau=0$. Gross value is $u[pv-(1-p)\ell s]$; all costs remain (\ref{eq:randomvalue}).

The finite global candidate list at zero width is $0,\tau_*,1$, feasible roots of
\begin{equation}
(T-1)+(2-T-\alpha)\tau-\tau^2=0,\qquad
1-T-(1-\alpha)\tau+\rho\tau^2=0,\label{eq:zero_boundaries}
\end{equation}
and the feasible clipped peak (\ref{eq:clippeak}) with $d_{\max}=1$, $B_0(-\ell)=0$. There are no additional interior candidates: the mixing frontier is strictly convex in $u$, and on the all-submit posterior-indifference branch
\begin{equation}
W'(\tau)=-a-(1-p)C\left[1-\alpha-\frac{1-T}{\tau}\right]\frac{1-T}{\tau^2}<0.\label{eq:zero_decreasing}
\end{equation}
Complete derivations, including zero-width support and attainment, appear in Appendix \ref{app:random_opt}. Thus no economically admissible shock width is excluded from the optimized uniform-cost game.

\subsection{Expanded governance menu and completed synthetic results}\label{subsec:random_results}
Let $O^*=\max_{\tau\in[0,1]}W(\tau)$. The enlarged policy menu has value
\begin{equation}\max\{V_\varnothing,O^*,W_{01},W_{11}\},\label{eq:newmenu}\end{equation}
where locked policies retain their original absent/full auditing choices and genuine clipping correspondence. Equivalently replace $\max\{V_\varnothing,O_0,O_1\}$ in (\ref{eq:kglobal}) by $\max\{V_\varnothing,O^*\}$. This optimizes random auditing under open metrics only, not a continuous locked-audit menu or arbitrary mechanism. The full-range R2 endpoint boundary remains valid: a valuable unaudited open architecture is not harmed by costless \emph{full} audit. Nothing here reinstates a disproved pure-substitution reversal.

All numbers below are synthetic. The central R3 review supplied the $y$, deterrence thresholds and boundary/full-audit values for the actual R2 witnesses A and B; those are reproduced checks, not author-origin discoveries. The finite candidate comparison now establishes their global intensity choices:
\begin{center}\small
\begin{tabular}{lrr}\toprule
Quantity & R2 A & R2 B\\\midrule
$\tau_*$ & .679371013286 & .736677947435\\
$W(0)$ & -.070000000000 & .001617525773\\
$W(\tau_*)$ & .119015376935 & -.015522826262\\
$W(1)$ & .109885313147 & -.036055708365\\
Global intensity & $\tau_*$ & $0$\\
Optimized no evidence & .090000000000 & .000000000000\\\bottomrule
\end{tabular}\end{center}
A has an additional interior all-submit stationary point at $\tau=.575321168133$, value $-.001054316163$, and a mixed/all-submit boundary at $.609503597812$, value $-.001412028654$. Both lose to $\tau_*$. B has two mixed/all-submit boundaries at $.453567401325$ and $.645845300804$, and an interior all-submit stationary point at $.470755519023$, value $-.037002297847$. Its global choice is zero. Both optimized open options beat all other items of their enlarged governance menus.

A further exact-primitive witness establishes that a deterrence-only recommendation is false. Set $p=.5$, $v=1.6$, $\ell=2$, $r=.1$, $f=.00525$, $C=.105$, $\kappa=e=.001$, $a=.005$, $K=1$ and $\delta=.02$. The unique global intensity is the upper-clipped interior candidate
\begin{equation}\tau_c=3969/8200=.484024390244,\quad W(\tau_c)=.024142149390.\label{eq:interiorwitness}\end{equation}
It has $M=.510975609756$, $I=.489024390244$, $s=.763300267698$, $Q=1$ and $u=1$. Inflation strictly beats the negative default. The passed-packet mean is about $.041626$, above the largest fallback draw $.02$; it does not reveal quality. The exact candidate set also includes clipping boundaries $.357191915763$ and $.640308084237$, mixed/all-submit boundaries $.196437732657$ and $.935419172704$, $\tau_*=.936570748362$, zero and one. The best other candidate value is $.008438684154$. No feasible interior stationary point of the remaining all-submit branch improves it. No evidence has value zero, while the locked alternatives have values $-1.004000737854$ and $-.204000000000$. The same clipped optimum survives at $\delta=0$, as separately solved and checked; this does not alter the old binary-menu result.

The interior optimum is not created solely by positive checking expenditure. Changing only $a$ to zero gives the uniquely optimal $\tau=399/820=.486585365854$, with $W=.026568673780$, whereas $W(0)=-.003935112854$ and $W(1)=-.005954403634$. The unaudited starting architecture is not valuable, so this partial-versus-full comparison does not contradict or reinstate the disproved claim addressed by Theorem \ref{thm:boundary}. It is a different, explicitly solved segment of the expanded frontier.

The package also records the full optimal metric-only intensity interval when $a=0$, solves a $\delta=5$ case with funding after failed audits and a $\delta=12$ case where inflation is absent even at zero intensity. Raw checks enumerate quality, action, actual checking and observed message before evaluating the history-dependent receiver decision. They charge every completed check, including a failure and a checked accurate packet. They separately verify all sender deviations and the correct $f+\tau rd$ cutoff. There is no fabricated firm dataset, larger portfolio simulation or completed field study.

\section{Sponsor information and blinding through an explicit experiment}\label{sec:sponsor}
This section concerns the zero-shock baseline only. No global-shock sponsor ranking is asserted. No fixed political premium appears in the model. To ask about sponsor visibility, introduce an observable sponsor class $S$ with known conditional prior $p_s=\Pr(G\mid S=s)$. Its signal is genuinely informative about project quality. Conditional on quality, the cost distribution $F$ and $r,f$ do not depend on $S$. Visibility reveals $S$ before reporting to both manager and evaluator; blinding hides it from the evaluator. For the blinded comparison, the manager may know $S$, but its payoff depends on quality, cost and the common funding probability only, so this extra knowledge does not change its actions. All conditional priors are assumed to satisfy region R.

Visibility is a different information partition, not an adjustable information-retention coefficient. Solve Theorems \ref{thm:eq}--\ref{thm:policy} within each visible class using $T_s=p_sv/[(1-p_s)\ell]$, and average the resulting costs and values. Under blinding, use the marginal $p=\E[p_s]$.
\begin{proposition}[An explicit sponsor-signal comparison]\label{prop:sponsor}
Let $c\sim U[0,C]$, and suppose region R holds for every visible prior and for their mean. With open metrics and full factual audit, revealing the informative sponsor class weakly lowers ex ante firm net value. The loss relative to blinding is
\begin{equation}\frac{Cv^2}{2\ell^2}\left[\E\!\left(\frac{p_s^2}{1-p_s}\right)-\frac{p^2}{1-p}\right],\label{eq:sponsor}\end{equation}
strictly positive for nonconstant $p_s$. Without an audit, or with locked metrics, the visible and blind values coincide in this parameter region. These statements condition on the given architecture, before selecting the optimal menu item.
\end{proposition}
Here blinding improves the open-audit architecture by reducing resource dissipation. It does not harm allocation by deleting an arbitrary fraction of context. This reverses the rhetorical role of the old blinding example. The result does not contradict the value of information for a fixed decision problem: managers change their private-cost thresholds when sponsor class becomes observable. With full audit and locked metrics, quality is already revealed and the sponsor signal has no further decision value. If a visible class falls outside R, the comparison must be re-solved; no global welfare ranking of anonymity follows.

\section{Completed synthetic equilibrium analysis}\label{sec:numerics}
Every number in this section is a synthetic illustration in arbitrary units. No firm observations, treatments or estimated parameters are used. Set $p=.4$, $v=\ell=2$, $r=1$, $f=.2$, $\bar c=1.2$, $\kappa=.02$, $e=.03$, $a=.04$ and $K=.06$. Uniform private selection cost gives $T=2/3$, $F(f)=1/6$ and $c_T=.8$.
\begin{table}[ht]\centering\small
\caption{Synthetic Bayesian equilibria under uniform private selection cost}
\begin{tabular}{lrrrrr}\toprule
Architecture & $u$ & Bad metric & Bad inflation & Prep. cost & Net value\\\midrule
Open, no audit & .200 & .16667 & .50000 & .07000 & $-.11400$\\
Open, full audit & .800 & .66667 & .00000 & .16000 & $-.23600$\\
Locked, no audit & .200 & .00000 & .66667 & .08000 & $-.18400$\\
Locked, full audit & 1.000 & .00000 & .00000 & .00000 & $.69200$\\\bottomrule
\end{tabular}
\end{table}
Bad action masses are conditional on bad quality, while preparation costs and net values are unconditional per registered opportunity. Positive-report good probability is $.5$ in the first three rows. Full audit under open metrics eliminates factual inflation but raises bad-project funding from $.08$ to $.32$ and good-project funding from $.08$ to $.32$. Gross value remains zero. The joint policy has value $.692$, compared with the no-acquisition fallback value zero. Its maximum affordable governance cost is $.752$ for these primitives; the active-only comparison gives the distinct threshold $.866$.

For a rising triangular density on $[0,1.2]$, $F(c)=(c/1.2)^2$, the equilibrium audit cutoff becomes $.9797959$. For a falling triangular density, $F(c)=2c/1.2-(c/1.2)^2$, it becomes $.5071797$. Both satisfy R. All type choices, aggregate inflation rates, posteriors and costs are recalculated under the new distribution. The audit-only net-value effects are respectively $-.2143900$ and $-.0632649$, compared with $-.1220000$ under uniform cost. This is equilibrium robustness over private-cost distributions, not the earlier fixed-action performance exercise under a triangular project-value density.

The R2 implementation solves all four uniform-cost policies over the full range, not merely the original three perturbations. Section \ref{subsec:witnesses} reports the new policy witnesses. The package independently checks the closed quadratic against bisection, sender deviations over private types, numerical receiver-payoff integration at all information sets, regime boundaries, the locked clipping correspondence and failed-packet costs. These are computational checks within assumed primitives, not empirical validation or a proof by simulation. The proof of Theorem \ref{thm:boundary} supplies the global sign restriction.

\section{Proposed empirical evidence: estimands and identification}\label{sec:empirical}
\subsection{Economic outcome before method}
No experiment or observational firm analysis has been conducted. The separate protocol defines a project-opportunity register and a budget-isolated cohort $c$, including withdrawals and rejected proposals. For a policy $\pi$, let its submitted set be $S_c(\pi)$, actual funding be $F_{ic}(\pi)$, scale $s_{ic}(\pi)$ and consistently valued net implemented project surplus be $V_{ic}(\pi,s)$. The primary economic target is
\begin{equation}\Delta=\E[W_c(\pi_1)-W_c(\pi_0)],\quad
W_c(\pi)=\sum_{i\in S_c(\pi)}F_{ic}(\pi)V_{ic}(\pi,s_{ic}(\pi))-C_c(\pi).\label{eq:empirical}\end{equation}
$C_c$ includes information production, preparation, packet assessment, factual checking and metric-governance costs once. Investment costs already inside $V$ are not charged again. Unspent capital receives a prespecified common opportunity return. If fallback opportunities vary across decisions, record their values under the same rule and include the realized fallback payoff on every unfunded decision. In the theoretical shock extension, reported welfare is incremental to always choosing that same fallback, and no-evidence allocation can still fund using its prior. It is not an observed zero-value control. Equity of funding across units is a separate outcome from economic value.

The directly observable proposed endpoint is 24-month discounted incremental net portfolio cash flow, less policy costs, relative to a frozen no-project baseline. This is narrower than lifetime risk-adjusted surplus. A common observable valuation kernel and complete follow-up are needed to interpret it as risk-adjusted value. Remaining project value must be followed longer or bounded; a sponsor terminal forecast does not point-identify welfare. Funded failures retain sunk costs, salvage and missing-outcome flags.
\subsection{Implemented policy value versus rejected labels}\label{subsec:labels}
With random assignment of isolated cohorts to implemented policies, complete capture of their funded-project outcomes and costs, and no cross-cohort spillovers, arm means identify the average implemented-policy difference. They do not require every rejected project's hypothetical funded return. Consistency maps an observed cohort portfolio to its assigned policy. Randomization balances potential policy outcomes across cohorts.

This does not justify ranking rejected proposals from historical funded-only labels. For an off-policy comparison on a fixed register, fixed project scales and additive policy-invariant potential funded values $Y_i$, define $d_i=\pi_1(i)-\pi_0(i)$. Observed values $O$ contribute $\sum_{i\in O}d_iY_i$. If missing values $M$ satisfy justified intervals $[L_i,U_i]$, the sharp rectangular bounds are
\begin{align}
\underline\Delta&=\sum_{i\in O}d_iY_i+\sum_{i\in M}\min\{d_iL_i,d_iU_i\}-\Delta C,\\
\overline\Delta&=\sum_{i\in O}d_iY_i+\sum_{i\in M}\max\{d_iL_i,d_iU_i\}-\Delta C.\label{eq:bounds}
\end{align}
The argument is elementary partial identification, not a new theorem. Cross-project synergies, policy-dependent implementation, variable scale or altered entry require a joint portfolio model instead of these formulas. The package retains the explicitly synthetic four-project example with bounds $[-6,6]$; its sign is unidentified.
\subsection{A test aligned with the revised mechanism}\label{subsec:design}
The core proposed trial now crosses independent factual verification with ex ante metric governance, rather than treating sponsor visibility as the main second factor. Register opportunities first, randomly assign budget-isolated cohorts to four architectures, and maintain comparable packet length and allotted review time. Metric governance must genuinely restrict an enforceable measurement universe; a stylistic template is insufficient. Record full source provenance in a protected research layer without feeding it selectively to committees across arms. A factorial interaction in net portfolio value is the direct policy contrast corresponding to (\ref{eq:interaction}), not proof of the model's unique mechanism.

Sponsor visibility is a separate, optional randomized first-stage study using a specified information partition and identical later contextual addenda. It identifies initial visibility, not political power. Predetermined sponsor influence measures rank and access, with expertise and implementation capacity as competing explanations. The model does not authorize identifying influence from numerical-score deviations.

Select all primary outcomes and multiplicity rules before assignment. The revised protocol declares the verification-by-governance interaction on net portfolio value as the primary confirmatory contrast, with main effects secondary and familywise adjustment specified. Baseline quality ratings are imperfect and not ground truth. Post-treatment metric changes, preparation effort and entry are outcomes or potential mediators, not total-effect controls.

If governance is assigned after register lock, the trial conditions on the opportunity pool. If announced before opportunity creation, it can change proposal entry and composition. These are different estimands. Persistent manager learning and common budgets may require randomization at a larger cluster than a single committee round. A small number of repeated decisions in one firm cannot be treated as many independent clusters. Full-register rejected projects are valuable for selection analysis even though their hypothetical returns remain unobserved.

No private employer data or administrative access has been inferred from affiliation. Public SEC financial data and Census MOPS do not, as verified, supply the required internal evidence recipes and rejected project register. Optional marginal funding lotteries require explicit further organizational authorization; none has been run or requested externally.

\section{Conclusion}\label{sec:limits}
The revised claim is an optimized hidden-action frontier with a preserved policy boundary. Random auditing changes both fraud survival and the privately optimal metric cutoff. Passed positives and all submitted packets must therefore be accounted for separately. The minimum intensity eliminating inflation is derived, not imposed; a globally optimal policy can lie below it and retain inflation. The complete frontier is piecewise rather than generally unimodal. Its finite candidates include both actual regime changes and all feasible stationary points. The inflation/negative mixing regime is strictly convex in its incremental incentive and cannot supply an interior maximum, while the all-submit regime can.

The full-range R2 result remains valid: if unaudited open evaluation is valuable, costless full audit cannot reduce its value in this uniform-cost technology. The positive-fee R2 example is retained as a restricted-menu calculation, not a rescue of pure-substitution harm. Enlarging the menu improves the A witness at the deterrence boundary, leaves B's optimal zero intensity unchanged, and permits a separate strict interior partial-inflation optimum. Full factual audit is dominated within open metrics when checking has a positive per-audit cost and lower intensity already deters inflation. Metric governance can still be optimal when its locked option beats the optimized open frontier and no evidence.

The mandatory good-record support excludes voluntary suppression equilibria. Hidden recipe provenance, rather than an irrational score, prevents full revelation. Metric locking is assumed perfectly enforceable and payoff-relevant; the paper does not solve how to enforce it in a complex firm. Private control benefits are fixed; project quality and initial information productivity are not endogenous. The one-candidate capacity constraint is not a multi-project capital market. Optimality applies only to five stated menu options, with equilibrium-selection qualifications at locked clipping points. The full-range endpoint boundary does not assert a ranking between arbitrary partial intensities. The new global intensity theorem is for uniform costs and the explicitly pooled checked-pass information structure. Revealing the audit flag to the allocator, nonuniform costs or continuously optimized locked auditing requires a new solution. Those are material limits, not completed robustness exercises.

The source reconciliation now rests on actual Section 2, equations (6)--(10), and Appendix A of Wulf's June 2009 prepublication text. Its mirror authorization and identity with the authenticated final remain unverified. The 2017 and distinct 2019 Boleslavsky--Carlin--Cotton primary mechanisms already contain Bayesian allocation and evaluation-induced incentives. Ba\'nka--Chudziak v1 supplies a finite false-claim coverage result, not a truthful-selection impossibility theorem. Our short partition observation states an assumption's implication; the substantive result remains the solved equilibrium displacement and information/resource frontier. These comparisons support a bounded result-level distinction, not broad priority, substantial originality certification or top-journal readiness. No empirical prevalence or institutional validity follows from synthetic examples.

The delivered analysis consists of solved games, integrated proofs and synthetic reproduction checks. The separate empirical protocol is prospective. There have been no field outcomes, recruitment, outside contacts, spending or publication submissions. Randomized implemented-policy value and off-policy rejected-project labels remain different estimands.


{
\small
\setlength{\parindent}{0pt}
\section*{References}\label{sec:refs}
\small
\begin{enumerate}\setlength\itemsep{5pt}
\item Ba\'nka, F., and J. A. Chudziak. 2026. ``More Claims, Less Evidence.'' arXiv:2610.05547v1, submitted October 4. \url{https://arxiv.org/html/2610.05547v1}. Primary HTML inspected, especially Section 3 and benchmark construction/discussion. Version 1, not an authenticated conference final.
\item Boleslavsky, R., B. I. Carlin, and C. Cotton. 2017. \emph{Competing for Capital: Auditing and Credibility in Financial Reporting}. NBER Working Paper 23273. \url{https://doi.org/10.3386/w23273}. Complete March 2017 NBER PDF retrieved through its public URL and inspected, especially Sections 2--3 and Propositions 1--3. It is a working paper, not a verified final journal article. Source: \url{https://www.nber.org/system/files/working_papers/w23273/w23273.pdf}.
\item Boleslavsky, R., B. I. Carlin, and C. S. Cotton. 2019. \emph{Disincentive Effects of Evaluation}. Author-hosted manuscript dated February 26, 2019. \url{https://www.raphaelboleslavsky.com/s/Disincentive_Effects_of_Evaluation.pdf}. Complete 19-page primary manuscript inspected, especially Section 3 and Propositions 2--4; also listed as Queen's Working Paper 1410. Separately titled; no identity with reference 1 or a final publication is assumed.
\item Farina, A., G. R. Fr\'echette, A. Ispano, A. Lizzeri, and J. Perego. 2026. ``The Selective Disclosure of Evidence: An Experiment.'' \emph{Review of Economic Studies}, advance article rdag076, published July 10, 2026; corrected and typeset July 29. \url{https://doi.org/10.1093/restud/rdag076}. Publisher record verified. Inspected full author manuscript dated October 16, 2025: \url{https://www.jperego.com/files/papers/FFLP/draft.pdf}. It merges and replaces earlier 2024 papers; final main text is not asserted identical.
\item Hartmann, M., and B. E. Wei\ss enberger. 2024. ``Decision-making in the capital budgeting context: effects of type of decision aid and increases in information load.'' \emph{Journal of Business Economics} 94:379--411. Online August 17, 2023. \url{https://doi.org/10.1007/s11573-023-01165-5}. Open text inspected in original research round.
\item Lakkaraju, H., J. Kleinberg, J. Leskovec, J. Ludwig, and S. Mullainathan. 2017. ``The Selective Labels Problem: Evaluating Algorithmic Predictions in the Presence of Unobservables.'' \emph{KDD 2017}. \url{https://doi.org/10.1145/3097983.3098066}. Author PDF inspected in original research round; the present policy identification argument is separately proved.
\item Milgrom, P., and J. Roberts. 1986. ``Relying on the Information of Interested Parties.'' \emph{RAND Journal of Economics} 17(1):18--32. \url{https://doi.org/10.2307/2555625}. Author Stanford archive scan retrieved; 1985 Cowles Discussion Paper 749 is a separate working version. Archive: \url{https://web.stanford.edu/~milgrom/publishedarticles/}.
\item Nikias, A. D., S. T. Schwartz, and R. A. Young. 2021. ``The effect of information transparency on capital budgeting with privately informed agents: a short research note.'' \emph{Journal of Management Control} 32:253--268. \url{https://doi.org/10.1007/s00187-020-00311-2}. Publisher abstract and methodological notes inspected in original round; full main text unavailable.
\item Stein, J. C. 1989. ``Efficient Capital Markets, Inefficient Firms: A Model of Myopic Corporate Behavior.'' \emph{Quarterly Journal of Economics} 104(4):655--669. November issue, as printed in the final PDF. \url{https://doi.org/10.2307/2937861}. Author-hosted final PDF inspected, especially Section II: \url{https://stein.scholars.harvard.edu/sites/g/files/omnuum5951/files/stein/files/qje-1989.pdf}.
\item Stein, J. C. 2002. ``Information Production and Capital Allocation: Decentralized versus Hierarchical Firms.'' \emph{Journal of Finance} 57(5):1891--1921. \url{https://doi.org/10.1111/0022-1082.00483}. Author-hosted final PDF inspected, including model, Proposition 1, hard-information treatment and discussion. Author source: \href{https://stein.scholars.harvard.edu/sites/g/files/omnuum5951/files/stein/files/hierarchies-jf-final.pdf}{Harvard final article PDF}.
\item Wulf, J. 2009. ``Influence and inefficiency in the internal capital market.'' \emph{Journal of Economic Behavior \& Organization} 72(1):305--321. \url{https://doi.org/10.1016/j.jebo.2009.07.001}. Final publisher preview verified at \url{https://www.sciencedirect.com/science/article/pii/S0167268109001802}. Authenticated final full text was not obtained. June 2009 prepublication Section 2, equations (6)--(10), and Appendix A were read through the public index at \url{https://studylib.net/doc/10476638/published-in-journal-of-economic-behavior-and-organization-}. Host authorization and final-version identity remain unverified; July 2000 and September 2005 records are separate.
\end{enumerate}

}
\clearpage
\appendix
\section{Complete proofs and equilibrium support}\label{app:proofs}
\subsection{Proof of Theorem \ref{thm:eq}}\label{app:eq}
Because good projects automatically report $+$, its probability is at least $p>0$. Under every policy the type $G$ record therefore gives a Bayesian contribution $p$ to $+$. Failed audits cannot come from $G$. Negative records must have positive probability in any equilibrium. Otherwise every bad type would choose a positive action. With no audit this pools every type on the positive packet, whose prior expected value is negative, so it is rejected; costly advocacy would then be dominated by the nonnegative payoff from the negative default. With a full audit, inflation is dominated because failure reveals bad quality and is rejected. If all bad types instead selected metrics, positive packets would again pool every type and be rejected, contradicting costly selection. With audit and lock, inflation is directly dominated. Thus some bad types preserve the negative record, its posterior is zero, and not funding is the unique optimum. Bad-manager payoffs are then (\ref{eq:payoffs}).

First consider $(A,L)=(0,0)$. If $u<f/r$, inflation gives negative payoff and is unused. Selection is strictly better than $N$ for $c<ru$ and strictly worse for $c>ru$, so $s=F(ru)<F(f)<T$. Equation (\ref{eq:posterior}) then gives $\nu>0$, requiring $u=1$, a contradiction. If $u>f/r$, inflation gives positive payoff, and every bad type produces a positive packet using its preferred instrument. Thus $s=1>T$, giving $\nu<0$ and requiring $u=0$, also a contradiction. Necessarily $u=f/r\in(0,1)$. Types $c<f$ strictly prefer selection since $ru-c=f-c>0$; $c>f$ obtain zero from inflation and $N$, and negative payoff from selection. For the receiver to mix, $s=T$. Hence inflation has mass $T-\alpha$, feasible because $0<T-\alpha<1-\alpha$. The stated constant mixing probability is one implementation. Every measurable assignment of inflation probabilities among $c>f$ with that total mass is another. All such assignments incur the same $f$ cost per inflating project. The type $c=f$ has zero mass and can take any of its tied actions. This verifies all deviations and establishes the complete aggregate correspondence.

For $(1,0)$, inflation always yields $-f<0$. Selection gives $ru-c$, so $s=F(ru)$. If $u=0$, $s=0$ and the receiver strictly funds the positive record, a contradiction. If $u=1$, $s=F(r)>T$ and the receiver strictly rejects, a contradiction. Therefore $u\in(0,1)$ and $F(ru)=T$. Strict monotonicity of $F$ gives $u=c_T/r\in(f/r,1)$. Selection versus negative is strict except at the measure-zero cutoff. There is no factual inflation on or off the manager's optimal strategy.

For $(0,1)$ the only bad actions are $I,N$. If $ru<f$, all preserve $N$, so $s=0$ and the receiver must set $u=1$, contradicting $ru<f$ since $r>f$. If $ru>f$, all inflate, so $s=1>T$ and the receiver must set $u=0$. Thus $u=f/r$ and the required aggregate inflation probability is $s=T$. Its allocation across private costs is unrestricted because those costs do not affect either feasible payoff in this policy. Both positive and negative events have positive support.

For $(1,1)$, inflation costs $f$ and cannot pass. All bad types strictly choose $N$. Positive records reveal good quality, so the allocator uniquely funds them; negative records reveal bad quality and are uniquely rejected. For the first three policies, substitution of $s=T$ in Bayes' rule gives $\mu=\ell/(v+\ell)$ and zero conditional expected project value, supporting exactly the displayed receiver mixing probabilities. A manager cannot generate an additional unmodeled free message. Audit failures under an audited architecture, if caused by a deviation to $I$, have posterior good probability zero by the evidence technology; the allocator rejects. At unaudited histories a failure is technologically impossible, so beliefs there do not affect any feasible deviation and can be set to zero. These statements specify all receiver information sets. Good-report support is technological; no skeptical off-path belief is used to discard a good positive report. This proves existence, completeness and the stated multiplicity.
\subsection{Proof of Corollary \ref{cor:regret}}
For an imperfect architecture, unconditional implemented value is $u[pv-(1-p)\ell T]=0$. The full-information policy funds $G$ only, giving $pv$. Its difference from the imperfect policy can also be written $pv(1-u)+(1-p)\ell Tu$, counting rejected good projects and funded bad projects. Since $(1-p)\ell T=pv$, the sum is $pv$. Substitution of the two equilibrium $u$ values gives the funding probabilities in the corollary. Under audit and lock, the full-information decision is implemented and regret is zero.
\subsection{Proof of Theorem \ref{thm:policy}}\label{app:policy}
Under open metrics without audit, low-cost selection contributes $I(f)$ conditional on a bad project, and inflators have aggregate mass $T-F(f)$ at cost $f$ each. The total is $J$. With full audit and open metrics, selection up to $c_T$ costs $I(c_T)$. Without audit and with locked metrics, a mass $T$ inflates at cost $f$, giving $fT$. Audit and lock leave zero strategic preparation. Multiply these conditional costs by $1-p$ once. The positive packet probability is $P$ in every imperfect architecture and $p$ under the perfect architecture. Gross value is zero in the former and $pv$ in the latter. Adding one initial cost $\kappa$, one packet evaluation cost, audit expenditure if chosen, and one governance cost if chosen gives (\ref{eq:w00})--(\ref{eq:w11}).

Because $F(c_T)=T$, subtraction gives
\[
I(c_T)-J=\int_f^{c_T}c\,dF(c)-f[T-F(f)]
=\int_f^{c_T}(c-f)\,dF(c)>0.
\]
Strictness uses $c_T>f$ and positive probability on that interval. This proves (\ref{eq:auditloss}), even when $a=0$. Likewise $fT-J=fF(f)-I(f)=D>0$, proving (\ref{eq:lockloss}). The interaction is obtained by direct subtraction:
\[
\Gamma=pv+(e+a)(P-p)+(1-p)[I(c_T)+fT-J].
\]
Use $P-p=(1-p)T$ and $fT-J=D$ for (\ref{eq:interaction}). All terms are nonnegative and $pv>0$.

The two strict dominance comparisons show that the active optimum is among $00,11$. Their difference is $pv-K+e(1-p)T-ap+(1-p)J$, proving (\ref{eq:active}), with a tie at equality. Every imperfect architecture has strictly negative net value: it has positive expected preparation cost under R, even if every administrative cost is zero. The no-acquisition policy avoids these costs. Its prior expected candidate value is $pv-(1-p)\ell<0$, so retaining capital is sequentially optimal without a record. Comparing zero with $W_{11}$ proves (\ref{eq:optimal}). The equality correspondence is explicit; no tie in governance is resolved without disclosure.
\subsection{Proof of Proposition \ref{prop:strict}}\label{app:strict}
Choose $\delta<\min\{v,\ell,(1-p)\ell-pv\}$, and small enough that $F(f)<H_\delta(u_0)<1$, and that under a fully positive bad pool the receiver's posterior value is below $-\delta$. These inequalities are feasible by R and continuity. Good quality then strictly warrants investment for all fallback draws once it is revealed, and bad quality never does. Without evidence the prior also never warrants investment for any fallback draw.

For no audit, suppose $u<u_0$. Bad selection gives $s=F(ru)<F(f)$. Its posterior value is bounded below by $\nu(F(f))>0$. Choose $\delta$ smaller also than this positive number. The receiver then funds for all fallback draws, contradicting $u<u_0<1$. If $u>u_0$, all bad types produce $+$, $s=1$, and the receiver rejects for all draws by the chosen upper bound on $\delta$, contradicting $u>u_0$. Thus $u=u_0$ and $\nu(s)=\delta(2u_0-1)$. Algebra solves this equation as $s=H_\delta(u_0)$. Its location in $(F(f),1)$ supports the stated mixtures. Removing the selection option leaves the same argument with $s=0$ below $u_0$, giving the locked case.

With audit and open metrics, $s=F(ru)$. Define the best-response function $B(u)=\Pr[t<\nu(F(ru))]$. It is continuous and nonincreasing; $B(0)=1$. For small $\delta$, $B(1)=0$ since $F(r)>T$. Hence the continuous strictly increasing function $u-B(u)$ has one zero in $(0,1)$. That zero must be in the receiver's interior response region, and thus satisfies (\ref{eq:strictroot}). Alternatively, the left side of that equation is strictly increasing, while
\[
H'_\delta(u)=-\frac{2\delta p(v+\ell)}{(1-p)[\ell+\delta(2u-1)]^2}<0.
\]
At $u_0$, $F(f)<H_\delta(u_0)$. At $u=1$, $H_\delta(1)<T<F(r)$. There is exactly one root in $(u_0,1)$. The manager's threshold is then $ru_A$. Audit and lock retain perfect revelation, so their outcome remains unchanged.

For an imperfect-policy equilibrium, the conditional optimal payoff relative to the fallback is
\[
\E_t[(\nu-t)_+]=\frac{(\nu+\delta)^2}{4\delta}=\delta u^2,
\]
where $\nu=\delta(2u-1)$ lies in $(-\delta,\delta)$. Multiplying by the positive packet probability proves (\ref{eq:gdelta}). A negative record never leads to candidate implementation, so it yields zero relative to its fallback. Expected value under full information is still $pv$ because $t$ is independent with mean zero and only good projects are implemented.

As $\delta\to0$, $s_0\to T$ and the root $u_A\to c_T/r$: every convergent subsequence solves $F(ru)=T$, whose solution is unique. The compact interval and continuity ensure convergence. All truncated preparation integrals, packet probabilities, evaluation costs and gross values are continuous in these arguments. Strict comparisons (\ref{eq:auditloss}), (\ref{eq:lockloss}), (\ref{eq:interaction}) and negativity of the three imperfect active policies therefore hold for some sufficiently small positive neighborhood. The perfect package and no-acquisition values are exactly unchanged under the chosen support of $t$, so their threshold is unchanged. Type-level and off-path statements follow as in Theorem \ref{thm:eq}. This supplies a complete local result, not a uniform bound over all primitive vectors or all $\delta$.
\subsection{Proof of Proposition \ref{prop:sponsor}}\label{app:sponsor}
Conditional on each visible class, replace $p$ by $p_s$ in the established equilibria. Under uniform cost, $I(x)=x^2/(2C)$ and $c_{T_s}=CT_s$. The full-audit open preparation cost is
\[
(1-p_s)I(CT_s)=\frac{Cv^2}{2\ell^2}\frac{p_s^2}{1-p_s}.
\]
Initial cost is fixed and the positive-packet probability $p_s+(1-p_s)T_s=p_s(1+v/\ell)$ is linear in $p_s$. Gross project value is zero in each class. Therefore all other open-audit costs have the same expectation under visibility and blinding. The sole difference is the displayed preparation term. The function $x^2/(1-x)$ has second derivative $2/(1-x)^3>0$ on $(0,1)$. Jensen's inequality proves (\ref{eq:sponsor}) and its strictness.

Without audit, $J(T_s)=I(f)+f[T_s-F(f)]$. Multiplication by $1-p_s$ gives an affine function of $p_s$, since $(1-p_s)T_s=p_sv/\ell$. Positive packet probabilities are affine too. In the locked/no-audit architecture, preparation is $(1-p_s)fT_s=fp_sv/\ell$, again affine. Under locked/full-audit, value is $p_s(v-e-a)-\kappa-K$, affine. Averaging thus preserves the other architecture values. Blinded managers' class knowledge does not affect the action inequalities because $p_s$ does not enter their payoff conditional on quality and common receiver probability. These arguments prove the complete comparison within the stated support.
\subsection{Proof of Theorem \ref{thm:global}}\label{app:global}
The receiver funds iff its conditional mean exceeds $t$, proving (\ref{eq:clip}). Integration of $(x-t)_+$ over the uniform shock gives (\ref{eq:option}); substituting the prior mean gives (\ref{eq:none}). Independence makes the posterior independent of the privately observed $t$. Negative and failed records can only come from bad quality, even off path, by the mandatory-record support and the audit technology. Hence their funding probability is $n$, and subtracting $rn$ from every sender payoff gives the stated action comparisons. Audited inflation is worse than $N$ by exactly $f>0$ for every type.

With audited open metrics, a bad manager selects iff $c<rD(s)$. Since $r<C$, uniform costs imply $s=\rho D(s)$. The left side is strictly increasing and the right side continuous nonincreasing. At zero the right side is positive, because $h_\delta(v)>h_\delta(-\ell)$. At $\rho$ it is strictly less than $\rho$ since $\nu(\rho)<0$ by R. Thus exactly one root $y\in(0,\rho)$ exists. The positive-packet funding response at this root is strictly interior. Its lower clipping would imply $D\leq0$ and contradict $y>0$. For upper clipping, if $\delta\leq\ell$ then $n=0$ and $y=\rho$, contradicting $\nu(\rho)<0$. If $\delta>\ell$, upper clipping requires $y=\rho(\delta+\ell)/(2\delta)$ and
\[
p(v+\ell)\geq(\delta+\ell)\left[p+(1-p)\rho(\delta+\ell)/(2\delta)\right].
\]
At $\delta=\ell$ this fails because $\nu(\rho)<0$. The right side is strictly increasing for $\delta\geq\ell$: its derivative is $p+(1-p)\rho(1-\ell^2/\delta^2)/2>0$. It therefore fails for every larger $\delta$ too. Consequently $u=n+y/\rho\in(0,1)$.

For $\delta\leq\ell$, $n=0$ and substituting $\nu(\rho d)=\delta(2d-1)$ gives
\[
2\delta(1-p)\rho d^2+[2\delta p+(\ell-\delta)(1-p)\rho]d-p(v+\delta)=0.
\]
For $\delta\geq\ell$, $n=(\delta-\ell)/(2\delta)$ and $D(s)=(\nu(s)+\ell)/(2\delta)$ in the interior. Thus
\[
2\delta(1-p)\rho d^2+2\delta p d-p(v+\ell)=0.
\]
Both have positive leading coefficient, positive constant on the right and exactly one positive root. Rationalizing the quadratic formula gives (\ref{eq:quadratic}), stable also near zero. Audit and lock remove selection and strictly deter inflation, so $s=0$ and the disclosed quality induces the two stated funding probabilities.

Without audit, if $D(s)>b$, inflation strictly dominates $N$; low-cost types $c<f$ select instead when feasible, while every remaining type inflates. Hence $s=1$ and this can be supported exactly when $D(1)>b$. If $D(s)<b$, inflation is unused. Open metrics then require the same fixed point $y$ as under audit, feasible below $\alpha$ exactly when $D(\alpha)<b$; locked metrics require $s=0$, feasible when $D(0)<b$. If $D(s)=b$, all feasible types below $f$ strictly prefer selection, and all other types are indifferent between inflation and $N$. This yields precisely the masses and support in the middle row. These three sender comparisons are exhaustive and also verify all deviations. Any mixture of high-cost types with the required total mass is feasible, and no cost is affected by its type-level implementation.

An interior receiver response solves $\nu(s)=\delta[2(n+b)-1]$, yielding (\ref{eq:inverse}). Since $b>0$, the target mean strictly exceeds $-\ell$, so the denominator is positive. When $n+b=1$, the receiver instead funds every passed positive and $\nu(s)\geq\delta$, which rearranges to (\ref{eq:plateau}). When $n+b>1$ no funding response can support the tie. For open metrics the clipping tie is impossible in R: $n+b=1$ requires $\delta>\ell$ and $b>1/2$, whereas
\[
\nu(\alpha)<\ell\frac{(1-p)\rho(1-b)}{p+(1-p)\rho b}<\ell(1-b)/b<\ell<\delta.
\]
The first inequality uses $pv<(1-p)\ell\rho$. For locked metrics the plateau can occur, for instance $p=.14,v=10,\ell=2,r=1,C=1.1,f=.75,\delta=4$: $n=.25$, $b=.75$, and every $s\in[0,.1627907]$ is supported. The region inequalities hold strictly. This also demonstrates that aggregate multiplicity cannot be dismissed as merely relabeling types. Good managers have no alternative feasible report; technologically impossible histories can be assigned arbitrary beliefs without affecting feasible deviations. This completes all receiver sets, support, existence and completeness.

\subsection{Proof of Proposition \ref{prop:menu}}\label{app:menu}
Conditional optimal payoff at a receiver information set with mean $x$ is $B_\delta(x)$ relative to the fallback. Multiplying by the probabilities of all three information sets gives (\ref{eq:allsets}). A selected lowest-cost mass $M$ costs $\int_0^{CM}c\,dc/C=CM^2/2$. Each inflation costs $f$; every submitted positive, including a failure, costs $e+Aa$. Add $\kappa$ and $LK$ once to obtain (\ref{eq:fullcost}). In equilibrium $AI=0$, giving (\ref{eq:gfull}). A smaller $s$ is more informative: obtain a larger-$s$ experiment from a smaller one by randomly relabeling some known-bad negatives as positives. Therefore optimizing the receiver's decision cannot decrease gross value as $s$ falls. This proves the gross-information-regret claim without using preparation costs as investment returns.

The finite-menu maximum is now explicit. Comparing the largest locked value before $K$ with the largest other value proves (\ref{eq:kglobal}), with all ties retained. Subtract $W_{00}$ from $W_{10}$, noting the common initial cost cancels, to obtain $P(y)(a^*-a)$ and (\ref{eq:afees}). On an upper-clipped interval $B_\delta(\nu(s))=\nu(s)$, so $P(s)B_\delta(\nu(s))=pv-(1-p)\ell s$. The negative contribution, inflation cost and packet cost are affine in $s$ too. Hence endpoint values give the exact range under aggregate multiplicity. They do not select an equilibrium on behalf of the firm.

\subsection{Proof of Theorem \ref{thm:boundary}}\label{app:boundary}
At $\delta=0$ all imperfect architectures have strictly negative value by Theorem \ref{thm:policy}, so the premise cannot hold. Consider $\delta>0$. A fully positive unaudited pool has gross value $B_\delta(m)=V_\varnothing$ and strictly positive preparation cost, so it too cannot satisfy the premise. If unaudited inflation is unused, the open audited and unaudited actions coincide at $y$, as do gross and preparation costs; at zero audit fee their values coincide. The open clipping plateau is impossible by Theorem \ref{thm:global}. It remains to consider genuine mixing. Write $x=s_0>\alpha$, $y=s_A$ and $q=p/(1-p)$. Since $D(x)=b$ and $y=\rho D(y)$, monotonicity gives $\alpha<y<x$. Equality of $y=x$ would require $x=\alpha$, excluded here. The posteriors of both passed packets are strictly inside the receiver's response interval. Gross value strictly increases when $x$ falls to $y$: the derivative in this interior is
\[
\mathcal G'(s)=(1-p)\{B_\delta(\nu(s))-(\nu(s)+\ell)h_\delta(\nu(s))-B_\delta(-\ell)\}<0.
\]
Specifically, with $u=h_\delta(\nu(s))$, the expression in braces is $u(\delta-\ell-\delta u)<0$ if $\delta\leq\ell$, and $-(\nu(s)+\ell)^2/(4\delta)<0$ if $\delta>\ell$.

First suppose $\delta\leq\ell$, so $n=0$. Let $G_0=\mathcal G(x)$ and $G_A=\mathcal G(y)$; let $C_0,C_A$ be preparation costs. Interior funding probabilities are $\alpha/\rho$ and $y/\rho$. Hence
\begin{align*}
G_0&=\delta(1-p)(q+x)\alpha^2/\rho^2,&
G_A&=\delta(1-p)(q+y)y^2/\rho^2,\\
C_0&=(1-p)C\alpha(x-\alpha/2),&
C_A&=(1-p)Cy^2/2.
\end{align*}
If $C_A\leq C_0$, strict gross improvement and reduced packet probability immediately give $W_{10}|_{a=0}>W_{00}$. Otherwise all quantities are positive, and
\[
\frac{G_A}{G_0}>\frac{C_A}{C_0}
\quad\Longleftrightarrow\quad
(q+y)(2x-\alpha)>(q+x)\alpha.
\]
The difference is $2q(x-\alpha)+2xy-\alpha x-\alpha y>0$, since $x>y>\alpha$. The hypothesis implies $G_0\geq C_0+\kappa+eP(x)+V_\varnothing\geq C_0$. Therefore
\[
G_A-G_0>G_0(C_A/C_0-1)\geq C_A-C_0.
\]
Adding saved packet costs $e[P(x)-P(y)]\geq0$ proves the strict net improvement.

Next suppose $\delta>\ell$. The interior identity is $D(s)=p(v+\ell)/(2\delta P(s))$. Equating $D(x)=\alpha/\rho$ and $D(y)=y/\rho$ gives
\[
\alpha(q+x)=y(q+y),\qquad
C_0-C_A=\frac{(1-p)C}{2}\{y^2-\alpha^2+2q(y-\alpha)\}>0.
\]
Audit both improves gross information and lowers preparation cost in this case, independently of whether the baseline is valuable. Evaluation cost weakly falls. Transition points follow by the same formulas or continuity. This proves the nonnegative threshold and strictness with inflation for every finite shock size.

The two fully specified witnesses in Section \ref{subsec:witnesses}, computed directly from (\ref{eq:quadratic})--(\ref{eq:fullcost}), have strict comparisons with every other menu item and no evidence. Their regime inequalities are strict. Thus their policy rankings persist under sufficiently small changes in primitives by continuity. Example B has $C_A>C_0$ but $a^*>0$; its negative actual effect requires $a>a^*$. This proves the existence statements and the boundary simultaneously. It does not extend the boundary theorem to nonuniform cost distributions or intermediate audit intensities.

\subsection{Proof of Theorem \ref{thm:random_eq}}\label{app:random_eq}
Conditional on bad quality, an inflator is funded with probability $(1-\tau)h_\delta(\nu(s))+\tau n$. Selection is funded with probability $h_\delta(\nu(s))$ and the unchanged negative with $n$. Subtracting $rn$ yields (\ref{eq:randompay}). Consider a proposed incremental probability $d>0$ and set $\gamma=(1-\tau)d-b$.

If $\gamma<0$, inflation is dominated by $N$ and exactly the types $c<rd$ select, giving $M=\rho d$, $I=0$, $s=\rho d$. If $\gamma>0$, inflation strictly beats $N$; comparing selection with inflation gives $rd-c>(1-\tau)rd-f$ exactly when $c<f+\tau rd$. This cutoff is below $rd$ since $\gamma>0$, hence below $r<C$. Therefore $M=\alpha+\rho\tau d$, $I=1-M$, $s=1-\tau+\tau\alpha+\rho\tau^2d$. If $\gamma=0$, the two cutoffs agree: $f+\tau rd=rd$. Types below $rd$ strictly select; higher types are indifferent between $I$ and $N$, giving every $s$ in $[\rho d,\rho d+(1-\tau)(1-\rho d)]$. These comparisons exhaust every feasible action and all sender deviations; cutoff types have probability zero.

This sender passed-mass correspondence is nondecreasing in $d$, with a filled upward jump at $d=b/(1-\tau)$. The jump is nonnegative because its length is $(1-\tau)(1-\rho d)>0$ whenever relevant. The receiver increment $D(s)$ is continuous nonincreasing. A metric-only outcome must equal the unique root $y=\rho D(y)$ already proved in Appendix \ref{app:global}, whose $d_y$ is positive and whose positive response is interior. Its inflation constraint is $(1-\tau)d_y\leq b$, equivalent to (\ref{eq:taustar}). If $d_y\leq b$, the metric-only equilibrium works at every intensity; otherwise $0<\tau_*<1$.

For $\tau<\tau_*$, $d_0<d_y$ implies $D(\rho d_0)>d_0$. If $D(U)\leq d_0$, continuity supplies an $s\in(\rho d_0,U]$ with $D(s)=d_0$. Since $d_0<d_y<1-n$, its positive funding probability is interior and the inverse is strictly decreasing there. Thus $s$ is unique and given by (\ref{eq:randommix}), with positive inflation and the stated feasible masses. If $D(U)>d_0$ and $\tau>0$, let $F_\tau(s)$ be the left side minus the right side of (\ref{eq:randomstrict}). This function is continuous and strictly increasing: its derivative where it exists is $1-\rho\tau^2D'(s)\geq1$. At $U$ its value is $\rho\tau^2[d_0-D(U)]<0$. At one its value is $\tau[1-\alpha-\rho\tau D(1)]>0$. To verify the last inequality, if $D(1)\geq d_0$ then $\alpha+\rho\tau D(1)\leq\rho D(1)\leq\rho<1$; if $D(1)<d_0$ then $\alpha+\rho\tau D(1)<\rho d_0<\rho<1$. Thus there is exactly one root $s\in(U,1)$. It must have $D(s)>d_0$: otherwise its right side is at most $U$, contradicting $s>U$. The strict sender cutoff is therefore below $rd<C$, and every proposed action is feasible. At $\tau=0$, the strict condition is $D(1)>b$ and (\ref{eq:randomstrict}) directly gives $s=1$, with the same verified sender actions.

Aggregate uniqueness follows from the opposite monotonicities except potentially where $D$ is flat at a jump. The lower flat region has $D=0$, which cannot equal $d_0>0$. The upper flat region has $D=1-n>d_y>d_0$, which also cannot equal it. Thus the jump always intersects the strictly decreasing part and fixes a unique mass. A strict branch with a flat receiver response instead fixes $s$ directly. No lower-clipped positive response is possible: if it were zero then $n=d=0$, and every costly action is dominated, implying $s=0$ and a strictly positive receiver response to the good packet, a contradiction. Upper clipping is allowed by the strict case.

At $\tau_*$, the metric-only equilibrium has $s=y$. Any aggregate inflation would increase $s$ above $y$, lower $D(s)$ locally strictly, and make inflation strictly worse than $N$. Hence it creates no additional equilibrium even though individual high-cost senders are indifferent at the boundary. A hypothetical distinct strict-branch equilibrium would contradict the monotonicity argument. Failed and negative messages are structurally bad because no good type can produce them. This determines their posterior even when they have zero probability on path. Passed positives always contain the good mass $p>0$. Scheduling flags are absent from the stipulated receiver information partition. At case boundaries, the formulas give the same $M,I,N,s$; continuity and uniqueness give continuity in intensity as well. Differentiability is not required at the boundaries. Finally, on $[\tau_*,1]$ only $a\tau Q$ varies, proving the strict full-audit dominance statement for $a>0$.

\subsection{Proof of Theorem \ref{thm:random_opt} and the zero-width extension}\label{app:random_opt}
On an interior positive response $\nu=2\delta(n+d)-\delta$. Since $\ell-\delta+2\delta n=L$, inversion of Bayes' rule yields $q+s=\lambda/H(d)$ and (\ref{eq:frontierS}). Also $B_\delta(-\ell)=\delta n^2$ for every positive $\delta$. Substituting the positive and bad-message option values into (\ref{eq:gfull}) gives (\ref{eq:frontierg}). Combining $I=(d/b)(s-\rho d)$ with the integral $CM^2/2$ gives (\ref{eq:mixfrontier})--(\ref{eq:mixW}). Feasibility of $I$ is exactly the displayed interval for $S(d)$.

We first exclude a mixed-regime interior maximum by direct curvature. Here $S'=-2\delta\lambda/H^2$, and $(dS)'=\lambda L/H^2-q$. Differentiating twice gives
\begin{align*}
g''(d)&=2(1-p)\lambda\delta L^2/H^3,\\
Q''(d)&=-\frac{1-p}{b}\{4\lambda\delta L/H^3+2\rho\},\\
[\tau(d)Q(d)]''&=-(1-p)\left\{\frac{2\rho}{b}+
\frac{4\lambda\delta(L+2\delta b)}{bH^3}+\frac{2bq}{d^3}\right\}.
\end{align*}
For $L=0$ the formula for $g''$ is zero, consistently with $g$ being affine. Preparation's negative has curvature $(1-p)r[4\lambda\delta L/H^3+\rho]$. Thus
\begin{align}
W_m''(d)={}&\frac{2(1-p)\lambda\delta L^2}{H^3}
 +(1-p)r\left(\frac{4\lambda\delta L}{H^3}+\rho\right)\notag\\
&+\frac{e(1-p)}b\left(\frac{4\lambda\delta L}{H^3}+2\rho\right)\notag\\
&+a(1-p)\left\{\frac{2\rho}b+
\frac{4\lambda\delta(L+2\delta b)}{bH^3}+\frac{2bq}{d^3}\right\}>0.\label{eq:mixconvex}
\end{align}
The strict inequality uses $r,\rho,1-p>0$ and nonnegative administrative costs. On every connected feasible mixed interval, an interior point has smaller value than at least one endpoint. The monotone reparameterization $\tau=1-b/d$ preserves this statement, although $W$ need not be convex in $\tau$ itself.

For a strict all-submit interval, $M=\alpha+\rho\tau d$. Expand $CM^2/2+f(1-M)$ to obtain $f-f^2/(2C)+C\rho^2\tau^2d^2/2$. Differentiating $S(d)=1-(1-\alpha)\tau+\rho\tau^2d$ gives (\ref{eq:dprime}); its denominator is positive. Differentiating $g(d)$ gives $g_d$ and hence (\ref{eq:strictderivative}). Multiplying by $H^2[\rho\tau^2+2\delta\lambda/H^2]>0$ gives precisely $R=0$. On the upper-clipped interval, $\mathcal G(s)=pv+(1-p)B_\delta(-\ell)-(1-p)[\ell+B_\delta(-\ell)]s$. Its substitution into (\ref{eq:randomvalue}) yields
\[
W_c'=(1-p)[\ell+B_\delta(-\ell)](1-\alpha)-a
 -(1-p)\{2[\ell+B_\delta(-\ell)]\rho d_{\max}+C\rho^2d_{\max}^2\}\tau.
\]
The slope is strictly negative, proving strict concavity and (\ref{eq:clippeak}). Only a feasible stationary point belongs to that branch.

We next establish finiteness and completeness of the procedure. A mixed/all-submit boundary has $N=0$ and $d=b/(1-\tau)$, so $S(d)=\rho d+b/d-b\rho$. Clearing the positive $dH$ gives the cubic (\ref{eq:transitionpoly}); its leading coefficient is $-2\delta\rho\ne0$. A clipping transition has $d=d_{\max}$, giving the stated quadratic. Thus finitely many actual boundary points partition $[0,\tau_*]$ into smooth regimes. At $\tau_*$ and at zero all possible changes are explicitly included. There are no other receiver or sender regimes by Theorem \ref{thm:random_eq}.

Interior strict stationary points can be found by a finite polynomial elimination. Write $V=q+1-(1-\alpha)\tau$. As a polynomial in $d$, $E$ has coefficients
\[
E_2=2\delta\rho\tau^2,\quad E_1=2\delta V+L\rho\tau^2,\quad E_0=LV-\lambda.
\]
For $\tau>0$, its derivative in $d$ is positive at the positive feasible root. Its discriminant is
\[
(2\delta V-L\rho\tau^2)^2+8\delta\rho\lambda\tau^2.
\]
It is not a square in the rational-function field in $\tau$. If $L>0$, comparing the leading, cubic, constant and linear coefficients of a putative quadratic square forces that square to be $(2\delta V-L\rho\tau^2)^2$, leaving the strictly positive additional quadratic term unmatched. If $L=0$, the expression is a quadratic with nonzero constant term and strictly negative discriminant as a polynomial in $\tau$, so it too is not a square. Hence $E$ is irreducible over this field. This observation also handles exceptional stationary identities rather than assuming a generic parameter vector.

If the resultant is nonzero, it has finitely many real roots. Every feasible solution of $E=R=0$ must be among them. Root filtering is necessary because algebraic conjugates, denominator clearing and nonfeasible actions can otherwise produce false candidates. A constructive elimination needs no general optimizer: pseudo-divide $R$ by the quadratic $E$ to a linear remainder $A(\tau)d+B(\tau)$, and form
\[
E_2 B^2-E_1AB+E_0 A^2.
\]
This differs from the resultant only by leading-coefficient factors irrelevant for $\tau>0$. Its real roots can be isolated, for rational inputs exactly by a Sturm sequence. If the resultant vanishes identically, irreducibility implies $E$ divides $R$. Then (\ref{eq:strictderivative}) vanishes identically on every feasible strict interior segment, and $W$ is constant there. Its already included closed-boundary values suffice, and the whole maximizing constant interval is retained. Thus this degeneracy does not create an infinite search problem.

A continuous equilibrium value on the compact intensity interval attains its maximum. Its maximizer is a boundary, an interior upper-clipped stationary point, an interior strict stationary point, or belongs to a constant segment. Mixed-regime interiors have already been excluded. The metric-only segment is affine with slope $-aQ(y)$, so its lower endpoint suffices unless it is constant, in which case all of it is retained if optimal. This proves necessity, sufficiency, finiteness of candidate values and the stated tie and attainment qualifications. Locked-policy multiplicity remains exactly as in Proposition \ref{prop:menu}; it is not removed by optimizing the open option.

At $\delta=0$, a negative or failed record is never funded. The receiver rejects positives for $s>T$, funds for $s<T$, and can mix at $s=T$. The same sender comparisons apply with $d=u$. For $\tau\geq1-\alpha/T$, $s=M=T$, $u=T/\rho$ satisfies every constraint. Below that threshold, $u_0=b/(1-\tau)<T/\rho<1$. If $T\leq U$, the filled sender jump contains $s=T$, yielding the stated mixing equilibrium. If $T>U$, the all-submit residual at $u_0$ lies below $T$, so funding must rise. Solving $s=T$ gives (\ref{eq:zero_u}); if its value exceeds one, the only supported outcome is $u=1$ with $s<T$. These are exhaustive and unique aggregates by the same monotonicity argument, now allowing the receiver's filled indifference interval. At $\tau=0$ R ensures the mixed branch; no division by $\tau$ occurs there. Cutoff and off-path arguments are unchanged.

The zero-width mixed frontier has $S=T$, $g=0$ and the same $\tau(d)$ and costs. Its second derivative is
\[
(1-p)r\rho+2(e+a)(1-p)\rho/b+2a(1-p)bq/d^3>0,
\]
so it has no interior maximum. The all-submit indifference branch has
\[
W= -\kappa-e-a\tau-(1-p)[f-f^2/(2C)]
 -(1-p)C[1-\alpha-(1-T)/\tau]^2/2.
\]
The bracket is positive because $u>0$, which gives (\ref{eq:zero_decreasing}). Its endpoints suffice. Clearing $T=U$ and $T=1-(1-\alpha)\tau+\rho\tau^2$ gives exactly the two quadratics (\ref{eq:zero_boundaries}). The strictly concave upper-clipped branch has the peak (\ref{eq:clippeak}) at $d_{\max}=1$, $B_0(-\ell)=0$. Compactness, boundary continuity and the same metric-only affine segment prove attainment and the complete finite list at zero width. This separately establishes the endpoint, rather than inferring it from an unjustified interchange of a limit and optimization.

\subsection{Verification of the strict interior witness}\label{app:random_witness}
The primitives in (\ref{eq:interiorwitness}) are rational, satisfy $0<\alpha=.05<T=.8<\rho=20/21<1$, and give $n=0$, $d_{\max}=1$. Substitution into (\ref{eq:clippeak}) gives exactly $3969/8200$. At this intensity $M=\alpha+\rho\tau_c$, $I=1-M$ and $s=1-\tau_c+\tau_c M$. They give $\nu(s)=.041626185823>\delta=.02$ and $(1-\tau_c)rd-f>0$, both strictly. Thus the upper-clipped branch and all sender actions are feasible. Direct application of (\ref{eq:randomvalue}) gives the displayed value, including actual audit checks on good, selected and inflating packets.

The two transition polynomials and the strict stationary polynomial are constructed with exact rational coefficients in the supplied code. Sturm isolation enumerates all their real roots in $(0,1)$, not just sign changes on a grid. Filtering leaves exactly the eight candidate points reported in Section \ref{subsec:random_results}; the output records the full stationary polynomial and every candidate value. Among them $\tau_c$ strictly beats the best alternative by more than $.0157$. The analytic candidate theorem proves that no unexamined intensity can do better. The differences with no evidence and both locked items are also strict, proving a uniquely optimal partial-inflation architecture in an open neighborhood of these positive-width primitives. At zero width the same peak satisfies its separately solved clipping constraints and candidate comparison. These are synthetic mathematical witnesses, not estimated economic magnitudes or experimental findings.

\section{Identification arguments}\label{app:identification}
For randomized implemented policies, assignment is independent of potential cohort portfolios. Consistency and no inter-cohort interference imply $\E[W_c\mid\pi_j]=\E[W_c(\pi_j)]$, when funded cash flows and all policy costs are fully measured. The difference of arm means identifies the average assigned-policy value contrast. This is an intent-to-treat estimand, including noncompliance and later unblinding where relevant. Rejected projects contribute no implemented cash flow under the realized policy; this never sets their potential funded value to zero. Missing funded outcomes still require follow-up or bounds.

For (\ref{eq:bounds}), minimize or maximize $d_iY_i$ over $Y_i\in[L_i,U_i]$ for each missing project. A positive coefficient is minimized at $L_i$ and a negative coefficient at $U_i$; the reverse gives maxima. Independent rectangular restrictions allow all endpoints to be attained jointly, proving sharpness. Known portfolio synergies or joint feasibility restrictions would replace the rectangle and its termwise solution. These logical bounds do not include sampling uncertainty unless additional confidence regions are constructed.

\end{document}
